\clearpage
\part*{Part III. Quantum Mechanics}
\label{atlaspart:03-quantum-mechanics}
\addcontentsline{toc}{part}{Part III. Quantum Mechanics}
\section*{Introduction to Part III}\label{narr:part:II}
\addcontentsline{toc}{section}{Introduction to Part III}
The framework of Part I speaks of pairs, links, reads and records, and never of a state vector. Part III asks how far the ordinary apparatus of quantum mechanics, the state space, the probability rule, the unitary step, the uncertainty relation, the measurement and the statistics of identical bodies, can be read off from that framework. The answer the Part gives is a sequence of conditional theorems, each resting on a short list of results from Part I, and a small number of open results where the reading parts. The order of the chapters is the order of the construction: the state space is built first from the outputs of the acts, the probability and the step are read at that space, the uncertainty relation follows from the step, measurement is the wash at one act, and the remaining chapters collect what the construction says about exchange, energy, the sum over records, the phase-space picture, the corollaries and the open system.

A reader coming from the textbooks will recognize the list, and the point of the Part is that it is not a list of postulates. Every item is read off objects that Part I has already built, the pair with its two ends, the link with its genre, the read in a basis and the records a read leaves, so that each chapter adds one reading of an object already in hand rather than one more assumption. The order matters because each reading uses the one ahead of it: without the state space there is no overlap to square, without the step there is no relation between two reads to bound, and without the wash there is no measurement to assign a statistics to. The chapters therefore run from the most basic object, the span of what the acts produce, to the most composite one, the open system that a sub-web presents to its complement. The conditional standing of the theorems is stated inside each box and is the same at every one, the axiom text and the chapters named as premises, each at its own standing. Where a chapter carries an open result, the box states both faces and names the parting, and this introduction names the open result beside the theorems so that the reader meets no surprise at the page.

Chapter 2-1 builds the state space. Theorem~35.1 states that the state space of a pair type is the span of the act outputs, so that the vector space the literature postulates appears here as the span of what the moves of chapter 1-4 produce. The same box reads the act at every one of the six pair types as a rank-one projection onto a pair, the span of its outputs under every admissible identification of bases as the whole register, the twirl over those identifications as a linear channel, and the inner product of the span as the Born readout, so that a qudit is its frame and superposition is the span of the frames. Three open results stand beside it, the Bell bound at every basis of the read clause, Open result~36.1, the algebra at the cap, Open result~37.1, and the attach-site basis read off the saturation waves, Open result~38.1, each with the point at which its derivation parts. The first of these is the one a reader of Bell will look for: at the pair state the four-term combination over the admitted bases reaches the standing bound at every one of the six pair types and stands above the local bound of two at every one, with the maximum over every basis decided as an exact element at the three types of modulus two and undecided at the three of modulus three and four, Open result~36.1. The second builds one algebra at the cap at every pair type, the direct sum of full matrix algebras on the record supports, and finds that the reversal of a hub leg is a symmetry of it at no hub-leg type, Open result~37.1. The third re-derives the saturation chain of every stable particle wave by wave and finds the hub count equal to the attach count at every one of the 745 lepton-form sinks, with a basis named by the saturation alone at 594 of them and the proton sink the first without one, Open result~38.1.

Chapter 2-2 reads the probability rule and the step at that space. Theorem~40.1 gives the flat pair with the unitary transport, the join and the step, which is where the count of records of chapter 1-3 becomes a square of an overlap. The mechanism is a single fact of linear algebra stated in the box: a pure pair on a common support has a flat reduced state exactly when its coefficient matrix is a multiple of a unitary, and the standing pair state is such a multiple at every alphabet dimension and every pair type, so the transport and the join are unitary as corollaries and the order step is unitary by its central generator, Theorem~40.1. The box checks this at the tabulated objects, the read transport a scaled partial isometry at all 26870 tabulated matchings, the join at a shared qudit a unitary at all 64 triples and the order step a unitary at the twelve steps, Theorem~40.1. Chapter 2-3 draws the uncertainty relation from the Weyl relation in Theorem~41.1. Two reads of the one link in two bases do not commute: the commutator of the two label observables is a nonzero exact matrix at every pair of distinct bases at the dimensions two, three and four and vanishes at a basis against itself, the variance of every read equals the norm of the centered read, and the product of two variances is bounded below by the squared modulus of half the commutator expectation at every pair of distinct bases and every register, with no violating pair, Theorem~41.1. The relation is therefore not imposed on the reads but read off the fact that the one link admits several bases.

Chapter 2-4 reads measurement as the wash at one act, with the collective readout of a joined web of several bodies in Theorem~42.1, the result the literature would file under the positive operator-valued measure. At every pairing the docks admit, the collective readout of the joined web of several bodies is defined at all six joined webs at two, three and four bodies, the interference terms between the copies exact and the coherence ratio an exact output at every one, Theorem~42.1. The box also names the reading beside at which the readout parts, the star at four bodies refused at the valence bound, so that the reader sees where the construction stops as well as where it reaches, Theorem~42.1. This chapter is the hinge of the Part, since the wash at one act is the object the chapters beyond it reuse: the exchange of bodies in chapter 2-5, the open system in chapter 2-10 and the projection postulate among the corollaries of chapter 2-9 are each a reading of it.

Chapters 2-5 to 2-8 carry the construction into the territory of identical bodies and energy. Theorem~43.1 obtains indistinguishability from the twirl, Theorem~44.1 reads energy as the first-power operator and the step, Open result~45.1 states, as an open result, the sum over paths as the finite sum over records, and Theorem~46.1 reads the phase-space function and the many-body sign. The thread through these four chapters is one idea: every quantity of the dynamics is a count or a sum over the records a read leaves, and the construction has only finitely many records to sum.

In chapter 2-5 the two copies of a site are taken as one qudit of the squared dimension, the standing two-copy configuration is built at all eighteen objects, the swap is washed by the twirl of the identification group and the readout is symmetric under it, Theorem~43.1. The sign under the swap is plus one at the amplitude level at every object, while the odd-cycle sign of the exchange triangle, minus one, holds beside and obstructs nothing, Theorem~43.1. In chapter 2-6 the first-power chain sum is an operator on the record space of every one of the 309 tabulated reads with the first-power sum of the web as its eigenvalue, it commutes with the step, and the step generator is re-met at every pair type with one level and conserved at all 32851 census moves, Theorem~44.1. The product generator carries the first-power sum as its level, and the Born square of the body level equals the standing rest energy at all thirteen bodies, so the Schrodinger form holds as a composition of two established objects rather than as a postulate, Theorem~44.1.

Chapter 2-7 is an open result and is read as one. The path integral as the finite record sum is stated at its five components and re-met at the established tables, the grading at the orders zero to three, the measure at 140 links of sixteen webs, the termination at order three and the record space of the one link at 172 letters with the support rank 319 equal to the Born count, yet the first entry of its keying is unkeyed with the eighteen entries beside keyed, and that entry is the parting, Open result~45.1. Chapter 2-8 then reads the phase-space picture. The negativity of the Wigner marginals beyond two copies is the negativity of the swap-face square of the direct sum, the two objects carry the same sign at every one of the five copy counts from one body to the parent bound of three, the first negative count is three at both, and the total of the signed distribution equals the swap-face square plus the closure read at all four webs, with a second route from the many-body links reproducing the parent grids cell by cell, Theorem~46.1. The reader who knows the Wigner function as a quasi-probability will find its negativity here as the sign of a finite count, read at the same place as the many-body sign of chapter 2-5.

Chapter 2-9 gathers the corollaries. Theorem~47.1 places the quantum-mechanics corollaries at their established roots, which is the point at which the reader can see which of the familiar statements rest on which theorems of this Part. The twenty-one entries of the remainder page of the glossary are each mapped to a root and typed from the kind of the map and the standing of the root, six of them derived as corollaries with a sentence, a proof sketch and an exact check, among them the Wigner and Wigner-Eckart theorems at the rotation group of the derived directions, the optical theorem at the unitarity of the flat pair, and the projection postulate and decoherence at measurement as the wash at one act, Theorem~47.1. Chapter 2-10 ends the Part with the open system: the table titled The chapter counted empty (in chapter~\ref{mod:II-10-decoherence-and-the-open-system}, among the instruments) records the chapter as counted empty, and the theorem titled The open system as the wash at a sub-web (chapter~\ref{mod:II-10-D536583-psyo-t1}, among the instruments) supplies the one derived result, the open system as the wash at a sub-web. At every tabulated sub-web carrying a shared record, 1753 of the 1755 sub-webs in 85 equivalence classes, the wash of every shared record at the complement is a twirl of the shared end, a completely positive, unital, trace-preserving and idempotent map, and the off-diagonal cells of the reduced state in the standing cap basis reach zero at the first record at every equivalence class, exact at a finite count and never asymptotic, the theorem titled The open system as the wash at a sub-web (chapter~\ref{mod:II-10-D536583-psyo-t1}, among the instruments). Decoherence is thus not a limit but a count, which is the form every result of this Part takes.

In the language of the literature, this Part names the Hilbert space, the Born rule, unitarity, the Weyl relation, the positive operator-valued measure, spin statistics, the Hamiltonian, the path integral, the Wigner function and decoherence. The correspondence is by name at every one of them. The boxes state what is established, which is the finite structure beneath each name.

The dependency lines of this Part point back almost entirely into Part I. The state space rests on the axiom chapter, the link, the records and the moves. The probability rule rests on the state space. The uncertainty relation rests on the probability rule, and the corollaries rest on all of them. Measurement rests on the records of chapter 1-3, the probe of chapter 1-6 and the step of chapter 2-2. The exchange of bodies rests on the link, the moves, the probe and the state space. The sum over records rests on the records, the moves, the saturation and the energy chapter, and the phase-space picture on the records, the measurement and the exchange. A reader who follows the derivation graph will see that Part III is the first branching of the thread beyond the framework and that the Parts on symmetry, gauge structure and electrodynamics draw on its results by number. The Part works with the finite structure beneath each name, and where the construction stops short of a textbook statement, at the every-basis maximum of the Bell combination, at the star of four bodies and at the first entry of the path-integral keying, the box states the parting.
\section{The state space from the act outputs}\label{mod:II-1-state-space-from-the-act-outputs}
\chdesig{3-1}{2-1}{II-1-state-space-from-the-act-outputs}{II-1-D541352-tfpk-t1}
\label{narr:open:2-1}The state space is the first object of this Part, and Theorem~35.1 states it as the span of the act outputs of a pair type, the vector space of the literature read off the moves of chapter 1-4. Beside it stand three open results, the Bell bound at every basis of the read clause, Open result~36.1, the algebra at the cap, Open result~37.1, and the attach-site basis from the saturation waves, Open result~38.1.
\begin{chconditional}[The state space of a pair type is the span of the act outputs]\label{box:II-1-state-space-from-the-act-outputs:1}
This result is conditional on the axiom text and on the results of the chapters named as its premises in the derivation below, each at its own standing.
At every one of the six pair types (the genres) the act is a rank-one projection onto a pair, the span of its outputs under every admissible basis identification is the whole register, the twirl over the identifications is a linear channel, and the inner product of the span is the Born readout: a qudit is its frame and superposition is the span of the frames.
\end{chconditional}
\chstanding{a conditional theorem}
\chdeps{Definitions 1.1 to 1.6, Theorems 2.1 to 2.3, Open results 2.4 and 2.5, Theorem 2.6, Open results 2.7 and 2.8 and Theorems 2.9 and 2.10, Theorems 9.1 to 9.7 and Open result 9.8, Theorems 18.1 to 18.4 and Open results 18.5 and 18.6 and Open results 241.1 and 241.2, Theorems 241.3 and 241.4 and Open results 241.5 to 241.7}{the theorems of Chapters 1-5, 3-1 to 3-3, 3-5, 3-9, 4-5, 7-1}
At every alphabet dimension the twirl of the whole identification group has exactly one pure fixed state, the maximally entangled pair on the common support, its amplitude the root of the Born readout.
The uniform washed state at every cut register of the construction is the marginal of the product of the pair states, so every traceless operator has expectation zero there.
A basis is a label passed along the pair: at a free link a read fixes no direction at any support dimension above one, and the reading under which the pair fixes its own directions forces the support to one at every one of the twenty-two free links of the construction.
The face of an arm in a readout is fixed by the passage of that readout alone, one register at every one of the 363 tabulated readouts.
Nothing crosses a corridor by the corridor alone: at every one of the 271 one-end placements of the 28 moves the far marginal of the pair state is unmoved when the act is read as a channel, and moves at 64 placements only when the outcome is kept, the classical conditioning a corridor uses.
Five open results stand here. The qudit label is the \rosetta{the frame class}{an equivalence class of reference frames} of the chain through it at 757 among the 1334 sites under one reading and at 958 under the other. The only-basis requirement is a reading, a pair fixing a direction of its qudit only at a one-dimensional common support. The frame class is the chain class at the 757 chain-faithful sites and the junction class at the six three-way cuts, the 571 hubs at neither. As the open set at a read it is defined at 96 among the 1334 sites under the standing definition of open and at 1163 under the larger-end reading. The bound-pair rule meets four of its five faces at the standing spelling and all five at the spelling beside.
\begin{chbridge}
In standard quantum mechanics the state space of a finite system is a Hilbert space, a maximally entangled pair the state whose reduced operators are flat, and the reduced state at one end of a pair is unmoved by any operation at the other end \cite{Nielsen2000}; the twirl over a group of local unitaries projects onto its commutant \cite{Horodecki2009}. The correspondence is by name: the compendium supplies a table of entries at their grades carrying the Standard Model names, a representation census and a kernel of conserved counts at chapter~\ref{mod:FRAMEWORK}, chapter~\ref{mod:I-5-saturation-in-one-figure-seeds-waves-termination-sinks}, chapter~\ref{mod:III-5-spin-and-helicity}, the chapter titled The standard-model table: the seventy-five entries with their status (chapter~\ref{mod:IV-10-sm-table-the-75-entries-with-their-status}, among the instruments), chapter~\ref{mod:VI-7-mixing-the-washed-overlap-and-the-ckm-block}, chapter~\ref{mod:IX-6-binding-as-the-pairing-records-the-deuteron-and-the-alpha}, the chapter titled The results table re-tallied (chapter~\ref{mod:XVI-1-the-results-table-re-tallied-at-the-cut}), chapter~\ref{mod:XIV-7-what-remains-the-ungraded-page} and the wash at a sub-web as the channel of this theory, the Born count as the dressed measure at the full web, the record algebra at a link and the corollaries re-met from the named clauses at chapter~\ref{mod:I-3-records-reads-and-the-born-count-the-gain-lemma}, chapter~\ref{mod:I-6-read-at-a-register-the-probe-and-the-measurement}, chapter~\ref{mod:II-2-born-rule-and-unitarity}, chapter~\ref{mod:II-3-uncertainty-relation-from-the-weyl-relation}, chapter~\ref{mod:II-9-qm-corollaries-the-fifteen-entries}, the chapter titled Decoherence and the open system (chapter~\ref{mod:II-10-decoherence-and-the-open-system}, among the instruments), chapter~\ref{mod:III-6-p-c-t-and-cpt}, chapter~\ref{mod:V-1-fine-structure-constant-at-the-one-link}, chapter~\ref{mod:VIII-2-body-alone-office-rule}, chapter~\ref{mod:XI-7-thermodynamic-remainder-and-the-arrow}, chapter~\ref{mod:XIV-7-what-remains-the-ungraded-page}, and the domain named as the formalism of the textbook corpus at its own axioms and the object named as a structure of the Hilbert-space formalism correspond to these objects by name.
\end{chbridge}
\begin{chproof}[census]
The premises of this derivation are the axiom and the six definitions of chapter~\ref{mod:FRAMEWORK}, the open result of chapter~\ref{mod:I-1-axiom-and-the-six-definitions}, the theorem of chapter~\ref{mod:I-2-link-and-its-six-genres}, the theorem of chapter~\ref{mod:I-3-records-reads-and-the-born-count-the-gain-lemma} and the theorem of chapter~\ref{mod:I-4-moves-in-one-table}. The chain of premises runs back to the axiom through those chapters, and the longest such chain in this compendium is twelve chapters long.
The acts are built at every pair type, the span of their outputs is computed by closure and by the Gram matrix, the twirl is built as a channel and checked cell by cell, and the Born square of the Gram matrix is compared with the readout at 30592 tested pairs; the commutants, passages and far marginals are computed at every tabulated case, each theorem conditional on the axiom text and on the open results it rests on.
By exhaustive enumeration of the six pair types of the tabulated genres, the act at each type is built as the rank-one projection onto the pair on the common support and its output tested for rank one, idempotence, trace one and one pair in and one pair out, with the squared amplitude re-met at every type. The orbit of the pair ket under the standing identification group (the monomial fourth-root group at each end, of order four, thirty-two, three hundred and eighty-four and six thousand one hundred and forty-four at the dimensions one to four) is spanned by closure and by its Gram matrix, equal to the whole tensor space at every type, the twirl over the identifications is a linear channel at every type, and the Born square of the Gram matrix equals the readout at every one of the 30592 tested pairs of rays. The smallest case is the two-qubit type, with the common support of rank two, the register dimension four, the act of rank one with the squared amplitude of its output one quarter, an orbit of thirty-two vectors spanning the four-dimensional space, a tight frame with the frame constant sixteen and eight rays, the Gram identity at all sixty-four pairs of rays, the two Born readouts at the pair one half and the register readout one quarter.
The counter-tests alter the objects singly and read what moves, a one-pair alteration moving the span and frame cells, a nonlinear readout the linearity cells, a non-flat pair the pair-ket norm, the trivial group the span cells and the planted text the deciding cells. The ten controls read the same at both engines and a second engine built apart reproduced every cell.
The machine-checked package named at the reproducibility appendix certifies the statements of this chapter under the definitions of this compendium and proves no theorem of the literature.
\chcert{30592}
\end{chproof}
\section{The Bell bound at every basis of the read clause}\label{mod:II-1-D535335-w-ebrc}
\begin{chwall}[The Bell bound at every basis of the read clause]\label{box:II-1-D535335-w-ebrc:1}
The read clause of the axiom text admits one read for each basis of the end register, and the four-term Bell combination at the pair state over the admitted bases reaches the standing correlation bound at every one of the six pair types (two root two at the (2,2), (2,3) and (2,4) types; twelve plus eight root three, over nine, at the (3,3) and (3,4) types; the four-cyclic surd at (4,4)), above the local bound two at every type, so the quantum structure is an output of the read clause and the record is local at no admitted read. The every-basis maximum is decided as an exact element at the three types of modulus two and undecided at the three of modulus three and four, where the attained cell lies below the algebraic bound four and no operator-norm form of the many-outcome inequality is established, the (3,3) type the first undecided.
Open is the every-basis maximum at three and four outcomes, a theorem at two outcomes alone, where the two-outcome cell of the Bell-readout result is re-met, with the derivation that would decide it recorded and not yet run.
\end{chwall}
\chstanding{an open result}
\chdeps{Definitions 1.1 to 1.6 and Theorem 35.1}{no theorem beyond it at this release}
\begin{chbridge}
The Bell inequality bounds the correlations of any local hidden-variable model at two outcomes by two; quantum mechanics reaches the Tsirelson bound of two root two at the maximally entangled pair; and for more outcomes per side the inequalities of Collins, Gisin, Linden, Massar and Popescu are violated by amounts whose maxima over all bases are known only in part \cite{Bell1964,Tsirelson1980,CGLMP2002}. The correspondence is by name.
\end{chbridge}
\begin{chproof}[exact part and finite part]
The premises of this derivation are the axiom and the six definitions of chapter~\ref{mod:FRAMEWORK} and the theorem of chapter~\ref{mod:II-1-state-space-from-the-act-outputs}. The chain of premises runs back to the axiom through those chapters, and the longest such chain in this compendium is twelve chapters long.
The admitted bases are printed at every end of every pair type, the four-term combination is maximised at the pair state over the admitted sub-families and compared with the standing correlation bound and with the local bound, and the maximum is decided at modulus two by the free-algebra identity (the square of the combination equal to four plus the product of the two commutators, with the norm bound root eight, the cut reads of a mixed type reduced to the projective reads of the support) and left undecided at the moduli three and four.
\textit{The exact part.}
From the three Weyl bases of the qubit end, by the four-term combination of the correlators of the pair state over the pairs of bases, we obtain the maximum two as an exact element of the ring of the eighth root of unity, the vector with first coordinate two and the others zero, attained with the labels free. Over the union of the named families the maximum is two root two, equal to the correlation bound and above the local bound two, and at the modulus two the every-basis maximum is decided by the free-algebra identity with the norm bound.
\begin{equation}
S_{\max}^{\mathrm{Weyl}}(2,2) = 2, \qquad S_{\max}^{\mathrm{union}}(2,2) = 2\sqrt{2}, \qquad S^{2} = 4\cdot 1 + [A_{1},A_{2}]\,[B_{1},B_{2}], \qquad \|S\| \le \sqrt{8} .
\end{equation}
\textit{The finite part.}
The procedure. For each of the six pair types and each of its two ends, list the admitted bases of the read clause that the named sub-families supply (the Weyl bases, the eigenbases of the Weyl letters, the offset frames, the Fourier frames at the offsets, and the adapted bases), each basis a complete orthonormal set in the ring, the duplicates at the union struck. For every choice of two bases at each end form the four-term combination of the correlators of the pair state (the two-outcome structure at the modulus two, the many-outcome structure at the moduli three and four), each correlator an exact element of the ring spelled over the surd basis, and take the maximum over the choices as an exact element. Compare that maximum with the standing correlation bound of the pair type and with the local bound two. At the modulus two decide the maximum over every basis of the end register by the free-algebra identity, the square of the combination equal to four plus the product of the two commutators, with the operator-norm bound root eight. The smallest instance is the two-qubit type, the admitted bases per end are the three Weyl bases and the four offset frames, every offset frame complete in the ring. Over the Weyl bases alone the maximum is two, the vector with first coordinate two and the others zero in the ring of the eighth root of unity, attained at the labels free. Over the union of the families the maximum is two root two, equal to the correlation bound, above the local bound two, and the every-basis maximum is decided at two root two by the identity.
The first undecided case is the type $(3,3)$, where the attained cell lies below the algebraic bound four and no operator-norm form of the many-outcome inequality is established, so the parting is named there.
The counter-tests alter the objects singly and read what moves. The first alters the continuum reading at the two-qubit type and moves fifty-one deciding keys, the first of them the continuum maximum at that type. The second alters the reading at the type $(3,3)$ and moves twenty-four keys. The four counter-tests stand at their pinned letters and the fifteen controls read the same at both engines. The result was reproduced by a second engine built apart from the first, the two runs equal at every deciding key.
The machine-checked package named at the reproducibility appendix certifies the statements of this chapter under the definitions of this compendium and proves no theorem of the literature.
\chcert{145}
\end{chproof}
\begin{figure}[h]\centering
\setlength{\unitlength}{1pt}
\begin{picture}(260,150)
\qbezier(86.0,75.0)(130.0,75.0)(174.0,75.0)
\put(130.0,83.0){\makebox(0,0){\scriptsize $(2,2)$}}
\put(77.0,75.0){\circle{14}}
\put(77.0,75.0){\makebox(0,0){\footnotesize $2$}}
\put(77.0,85.0){\makebox(0,0)[b]{\scriptsize $A$}}
\put(183.0,75.0){\circle{14}}
\put(183.0,75.0){\makebox(0,0){\footnotesize $2$}}
\put(183.0,85.0){\makebox(0,0)[b]{\scriptsize $B$}}
\end{picture}
\caption{The smallest instance of this chapter, the two-qubit link with its two ends, one read at each end in a basis the read clause admits. The four-term combination is formed from two bases at each end.}
\end{figure}
\section{The algebra at the cap}\label{mod:II-1-D535787-w-alcp}
\begin{chwall}[The algebra at the cap]\label{box:II-1-D535787-w-alcp:1}
At every one of the six pair types one algebra at the cap is built: the direct sum, over the records at the cap, of the full matrix algebras on the record supports, the records its minimal central projections, with the record counts 57, 172, 417, 412, 797 and 1455, the registers 96, 319, 906, 1036, 2053 and 5100 equal to the Born counts, and the dimensions 192, 757, 2766, 4252, 9061 and 38580. The algebra carries the Weyl basis block by block, the graded poset of the passed ends and the transport as the exponential of the order operator at every type, but the reversal of a hub leg is a symmetry of it at no hub-leg type: at the (2,4) type the far sector of 90 records at dimension two parts from the centre sector of 126 records at dimension four under the swap, so the orientation of a hub leg is a cell of the record spectrum.
Open is the hub-leg reversal as a symmetry of the algebra, the orientation of a hub leg a cell of the record spectrum; whether the first power (the energy operator of chapter~\ref{mod:II-6-energy-and-the-hamiltonian-the-first-power-operator-and-the-step}) is central in the algebra is the matter of that chapter.
\end{chwall}
\chstanding{an open result}
\chdeps{Definitions 1.1 to 1.6 and Theorem 35.1}{no theorem beyond it at this release}
\begin{chbridge}
The observables of a finite quantum system form a direct sum of full matrix algebras, one block per superselection sector, with the Weyl operators a basis of each block and the uniqueness theorem of Stone and von Neumann fixing the representation of the Weyl relations up to equivalence \cite{vonNeumann1931,Weyl1928}. The correspondence is by name: the compendium supplies a step generator at the order coordinate, a first-power operator, a direction functional and an office at the common support carrying the energy, momentum and mass names at chapter~\ref{mod:I-3-records-reads-and-the-born-count-the-gain-lemma}, chapter~\ref{mod:I-6-read-at-a-register-the-probe-and-the-measurement}, chapter~\ref{mod:II-6-energy-and-the-hamiltonian-the-first-power-operator-and-the-step}, chapter~\ref{mod:III-1-order-coordinate-and-its-theorem}, chapter~\ref{mod:III-6-p-c-t-and-cpt}, chapter~\ref{mod:IV-6-conservation-as-the-kernel-of-the-move-matrix}, chapter~\ref{mod:V-7-running-of-the-coupling}, the chapter titled The Dirac equation at the dock (chapter~\ref{mod:V-9-dirac-equation-at-the-dock}, among the instruments), chapter~\ref{mod:VIII-1-mass-unit-and-the-first-power}, chapter~\ref{mod:VIII-5-unified-mass-readout}, chapter~\ref{mod:X-2-field-equation-at-the-counting-face}, the chapter titled The gravitational and cosmological constants table (chapter~\ref{mod:X-9-gravitational-and-cosmological-constants-table}, among the instruments), chapter~\ref{mod:XI-1-second-law-as-the-boundary-count-under-acts} and chapter~\ref{mod:XIV-7-what-remains-the-ungraded-page}, and the energy named as the eigenvalue of a Hamiltonian generating a continuous evolution and the Hamiltonian named generating a continuous evolution on a Hilbert space correspond to these objects by name.
\end{chbridge}
\begin{chproof}[finite computation]
The premises of this derivation are the axiom and the six definitions of chapter~\ref{mod:FRAMEWORK} and the theorem of chapter~\ref{mod:II-1-state-space-from-the-act-outputs}. The chain of premises runs back to the axiom through those chapters, and the longest such chain in this compendium is twelve chapters long.
The record spectrum of the gain lemma and the register of the cap are read at every pair type, the direct sum is built with its centre the records and its dimension the sum of the squares, the Weyl basis is laid block by block, the graded poset of the passed ends and the transport as the exponential of the order operator are checked at every type, and the hub-leg reversal is tested as a symmetry of the algebra at the hub-leg types, where it carries the far blocks to centre blocks of a different count and dimension; the four open results of this Part that ask for one object are answered at the algebra and part at the reversal.
The procedure. For each of the six pair types read the record spectrum of the gain lemma at the cap, the sectors with their type, the count of records in each and the dimension of each record support. Build the direct sum over the records of the full matrix algebras on the record supports. Its center is spanned by the records, its register is the sum over the records of the support dimensions, and its dimension is the sum of the squares of the support dimensions. Lay the Weyl basis block by block, check the graded poset of the passed ends (one rank per cover, the ratios fixed per kind and increasing, the diamond closed) and the transport as the exponential of the order operator, and test the reversal of a hub leg as a symmetry of the algebra at the hub-leg types by carrying each far block to the block the swap sends it to. The smallest instance is the two-qubit type, the four sectors hold twenty-four odd records of dimension one, eighteen mixed-even records of dimension two, twelve same-even records of dimension two and three closure records of dimension four, so the record count, the register and the dimension are
\begin{equation} \adjustbox{max width=\dimexpr\linewidth-3em\relax}{$\displaystyle \begin{gathered} N = 24 + 18 + 12 + 3 = 57, \\ D = 24\cdot 1 + 18\cdot 2 + 12\cdot 2 + 3\cdot 4 = 96, \\ \dim = 24\cdot 1^{2} + 18\cdot 2^{2} + 12\cdot 2^{2} + 3\cdot 4^{2} = 192 , \end{gathered}$} \end{equation} the register ninety-six equal to the Born count of the type and normalized at the stop condition, and the graded poset of the four sectors having ranks zero, one, one and two with both ratios two. The first parting is the type $(2,4)$, where the far sector holds ninety records of dimension two and the center sector one hundred and twenty-six records of dimension four, so the swap of the hub leg carries blocks to blocks of a different count and dimension and is a symmetry at no hub-leg type.
The counter-tests alter the objects singly and read what moves. A wrong grading (the sectors ranked by their index in place of the passed-end count) moves the first parting and the grading readings. The reversal read as the identity in place of the exchange of the far and center roles moves the symmetry cells and the first parting. The four counter-tests stand at their pinned letters and the fifteen controls read the same at both engines. The result was reproduced by a second engine built apart from the first, the two runs equal at every key.
The machine-checked package named at the reproducibility appendix certifies the statements of this chapter under the definitions of this compendium and proves no theorem of the literature.
\chcert{3310}
\end{chproof}
\section{The attach-site basis from the saturation waves}\label{mod:II-1-D538971-w-atsx}
\begin{chwall}[The attach-site basis from the saturation waves]\label{box:II-1-D538971-w-atsx:1}
With the saturation chain of every sink (every stable particle of the network) re-derived wave by wave from the tabulated parent and move columns, the hub count equals the attach count at every one of the 745 lepton-form sinks, every hub a root event or a hub-step unit of the saturation, so the disagreement of the two readings at the preceding result dissolves, and the three lepton labels of the comparison table are present among the attach counts, the generation count three re-met. The saturation alone names a basis at 594 of these 745 sinks: the first without one is the proton sink 3027, whose root hub attaches at the three qutrits of the junction, with 96 tau-sink sinks at a repeated attach site, 48 hub-free sinks at an empty event list and one at the bound-respecting root hub at three sites.
Open is the attach-site basis at the sinks the saturation does not name; the generation count three at both readings stands at chapter~\ref{mod:IV-3-representations-from-the-open-end-census}, the Standard Model list re-met at chapter~\ref{mod:I-5-saturation-in-one-figure-seeds-waves-termination-sinks}, and the three generations as attach sites, the companion derivation, at chapter~\ref{mod:VIII-4-three-generations-as-attach-sites}.
\end{chwall}
\chstanding{an open result}
\chdeps{Definitions 1.1 to 1.6, Theorem 35.1 and Open results 245.1 and 245.2, Theorem 245.3, Open results 245.4 to 245.8 and Theorem 245.9}{no theorem beyond it at this release}
\begin{chbridge}
The state space of a quantum system is a Hilbert space spanned by a basis of outcomes, and the state space of a composite is the tensor product of the spaces of its parts; a graph state is specified by its graph, the vertices the qudits and the edges the entangling acts, and the sequence of acts that builds it fixes its stabilizer basis up to the ordering of the acts \cite{Dirac1930,Hein2004}. The correspondence is by name.
\end{chbridge}
\begin{chproof}[census]
The premises of this derivation are the axiom and the six definitions of chapter~\ref{mod:FRAMEWORK}, the theorem of chapter~\ref{mod:II-1-state-space-from-the-act-outputs} and the theorem of chapter~\ref{mod:I-5-saturation-in-one-figure-seeds-waves-termination-sinks}. The chain of premises runs back to the axiom through those chapters, and the longest such chain in this compendium is twelve chapters long.
The saturation chain of every sink is re-derived wave by wave from the established columns with the origin of every site carried into its labels, no edge unresolved at either chain; the hub count and the attach count are compared sink by sink; the attach counts are read against the lepton labels of the comparison table; and the attach-site basis is read off the waves wherever the saturation names one, the sinks at which it names none listed with their kind.
By exhaustive enumeration of the 745 lepton-form sinks of the tabulated saturation (every sink of the network whose form is the lepton form, read from the parent and move columns of the established sink tables), the saturation chain of each sink is re-derived wave by wave with the origin of every site carried into its labels and no edge left unresolved at either chain. The property tested at every sink is that its hub count equals its attach count, the attach count being the number of root events and hub-step units of the saturation, and the two counts are equal at every one of the 745 sinks. The attach counts are then read against the three lepton labels of the comparison table, all three present, and the attach-site basis is read off the waves wherever the saturation names one, a basis at 594 among the 745 sinks and none at the remaining 151, which are listed with their kind by the three kinds the basis tables of this result count, 96 tau-sink sinks at a repeated attach site, 48 hub-free sinks at an empty event list, and seven sinks at a root event attaching at several sites at once, the six proton sinks and the one tau-sink sink at the bound-respecting root hub at three sites, so that 96 plus 48 plus 7 is 151 and 594 plus 151 is 745. The first sink at which the saturation names no basis is the proton sink numbered 3027, the smallest case worked here, it has eleven sites and thirteen links at depth three, one hub, and one root event at its tenth site attaching at three sites at once (the sixth, seventh and eighth), so its attach count is one, equal to its hub count, and its root event names no single attach site. The next sink numbered 3616 has fourteen sites, seventeen links, two hubs, a root event at three sites and a hub-step event at one site, its attach count two.
The counter-tests alter the objects singly and read what moves. Striking the root events moves thirty-one summary cells, the first of them the count cell. Striking the hub-step events moves nineteen. The three counter-tests stand at their pinned letters, the ten controls read the same at both engines, and the base re-run is equal at every one of its four captures. The result was reproduced by a second engine built apart from the first, the two runs equal at every key.
\chcert{745}
\end{chproof}
\section{The three-four pair ket at the hub}\label{mod:II-1-D541352-tfpk-t1}
\begin{chtheorem}[The three-four pair ket at the hub]\label{box:II-1-D541352-tfpk-t1:1}
The office of the (3,4) pair ket at the hub under the body-alone rule is a recorded cell with the move classes re-met and the leg unit twelve at its first power: the (3,4) pair ket, a qutrit site and a ququart site at one (3,4) link, the hub-to-core in-leg, is a word of two sites and one link, its first power twelve at the adopted reading, no hub, no dock set, the dressing one and the act the measure twelve, equal to the leg unit at five of five recorded cells, the pair part one hundred forty-four equal to the first power squared, the one dark link at the (3,4) law with the second-order read $\tfrac{137}{414720} + \tfrac{37}{259200}\sqrt{3}$ and the closure $\tfrac{732073}{29859840000} + \tfrac{5069}{373248000}\sqrt{3}$, the office $\tfrac{4299817692073}{29859840000} + \tfrac{5069}{373248000}\sqrt{3}$ exact and nonzero, equal by exact field equality to the free-genre office of the Higgs and carrier row and located at five of five parents, its genre at the order two of the readability row, its move classes the hub-forming move, the cut and the hub step re-met three of three at the forces row and equal to the list of the twelve-slot table; the office below sixteen, equal to two and above two of the twenty recorded offices; the first parting none.
\end{chtheorem}
\chstanding{theorem}
\chdeps{Definitions 1.1 to 1.6, Theorems 9.1 to 9.7 and Open result 9.8, Theorem 35.1, Theorems 185.1 to 185.12, Theorems 200.1 to 200.3 and Open result 200.4 and Open results 241.1 and 241.2, Theorems 241.3 and 241.4 and Open results 241.5 to 241.7}{no theorem beyond it at this release}
The (3,4) pair ket is located at the Higgs and carrier row, the unpassed in-legs of the five bodies, the Higgs-office unit, the twelve-slot table and the roster; its office is read by the body-alone rule at the adopted reading with the pair part, the closure and the sign, placed among the twenty recorded offices by the exact sign of the difference, and its read order and move classes are re-met at the readability and forces rows.

\begin{chbridge}
The Yukawa coupling of a fermion to the Higgs field fixes the mass of the fermion as the product of the coupling and the vacuum expectation value, one term per fermion in the Lagrangian \cite{Yukawa1935,Weinberg1967}. The derivation supplies the office of the (3,4) pair ket here, and the coupling is named by correspondence.
\end{chbridge}
\begin{chproof}[finite and exact]
The premises are the Higgs and carrier offices with the body-alone rule (chapter 6-10, a theorem), the unpassed in-legs at the five bodies (chapter 8-1, a theorem), the Higgs leg unit (chapter 6-10, an open result), the readability order at the six genres (chapter 1-2, a theorem), the forces as the move classes (chapter 1-4, an open result) and the twelve-slot table (chapter 4-8b, an instrument). The chain back to the axiom runs through chapter 2-1 and the framework chapter. The procedure. The input is the (3,4) pair ket, a qutrit site and a ququart site at one link, the register whole one hundred forty-four and cut one hundred twenty, the record space with the sectors three hundred seventy-three, two hundred forty, one hundred fifty-six and twenty-eight at seven hundred ninety-seven records. The enumeration rule locates the object at the five parents and the roster, reads its office under the body-alone rule at the adopted reading, the measure struck, the slot, the in-leg at (3,4), the dock rule and the act read, and compares the office by exact field equality with the twenty recorded offices; the test at each parent is location and at each office the sign of the difference. From the one link at the measure twelve with no hub, no dock set and the dressing one, by the first-power rule, we obtain
\begin{equation}
p = 12 , \qquad p^{2} = 144 ,
\end{equation}
the first power equal to the leg unit at five of five recorded cells and the pair part the first power squared. From the link, dark at the dock rule as a genre other than the dock genre, by the second-order read at the (3,4) law,
\begin{equation}
G_{2} = \tfrac{137}{414720} + \tfrac{37}{259200}\sqrt{3} , \qquad c = (p\, G_{2})^{2} = \tfrac{732073}{29859840000} + \tfrac{5069}{373248000}\sqrt{3} ,
\end{equation}
over twenty-eight closure records, and from the two, by the one rule, the office
\begin{equation}
\omega = p^{2} + c = \tfrac{4299817692073}{29859840000} + \tfrac{5069}{373248000}\sqrt{3} ,
\end{equation}
exact, positive and nonzero, equal to the free-genre office of the Higgs and carrier row and to its leg-piece read, each one (3,4) link at the first power twelve, below sixteen and above two of the twenty recorded offices, the (2,4) free genre and the (2,2) pair ket at zero. From the readability row at its three faces we obtain the read order two, dark, equal to the twelve-slot table at its unread slot, and from the forces row the move classes at (3,4), the hub-forming move of the electromagnetic force at sixteen edges, the cut and the hub step of the weak force at two thousand six hundred seventy-two and three thousand three hundred eight edges, equal to the list of the twelve-slot table. The smallest instance on the page is the object itself, two sites and one link, with the pair within the bodies at their in-leg lines beside, the electron at the act thirty-six at the dressing three, the muon at one hundred eight at the dressing nine and the tau and the proton at three hundred twenty-four at the dressing twenty-seven, every line the measure twelve multiplied by the dressing of the body.

The counter-tests put the pair ket at the dock genre, the ququart struck to a qubit, which moved the first power to zero and the office to zero. A second struck the leg unit, the qutrit struck to a ququart, which moved the first power to sixteen off the leg unit at twenty-five cells. A third struck a move class, the cut, which moved the set off the list at eight cells. A fourth struck the text at its gate, which left every cell struck, and a fifth replaced the object by the Higgs word of the roster, located at no parent and off the leg unit, which moved the deciding cell. Every control read the same at both engines, eighteen of eighteen, and the office was reproduced at a second engine, equal at every cell.
\chcert{797}
\end{chproof}
\begin{figure}[h]\centering
\begin{adjustbox}{max width=\linewidth}\begin{tabular}{llrrrr}\toprule
object & genre & measure & act & dressing & first power \\ \midrule
the qutrit site & & & & & \\
the ququart site & & & & & \\
the one link between them & (3,4) & 12 & 12 & 1 & 12 \\ \bottomrule
\end{tabular}\end{adjustbox}
\caption{The (3,4) pair ket as a word of two sites and one link at the adopted reading, with its measure, act, dressing and first power, the leg unit twelve.}
\end{figure}
\par\medskip\noindent{\itshape With the state space spanned, the next chapter reads the probability rule and the step at it.}\label{narr:close:2-1}\par
\section{The Born rule and unitarity}\label{mod:II-2-born-rule-and-unitarity}
\chdesig{3-2}{2-2}{II-2-born-rule-and-unitarity}{II-2-born-rule-and-unitarity}
\label{narr:open:2-2}This chapter reads the probability rule and the unitary step at the state space of chapter 2-1. Theorem~40.1 gives the flat pair with its unitary transport, the join and the step, the point at which the count of records of chapter 1-3 becomes the square of an overlap.
\begin{chconditional}[The flat pair and the unitary transport, join and step]\label{box:II-2-born-rule-and-unitarity:1}
This result is conditional on the axiom text and on the results of the chapters named as its premises in the derivation below, each at its own standing.
A pure pair on a common support of rank $m$ has a flat reduced state exactly when its coefficient matrix is a multiple of a unitary on the support, and the standing pair state is such a multiple at every alphabet dimension and every pair type, so the transport and the join are unitary as corollaries and the order step is unitary by its central generator: the read transport is a scaled partial isometry at all 26870 tabulated matchings, the join at a shared qudit a unitary at all 64 triples, and the order step a unitary at the twelve steps.
\end{chconditional}
\chstanding{a conditional theorem}
\chdeps{Definitions 1.1 to 1.6, Theorems 2.1 to 2.3, Open results 2.4 and 2.5, Theorem 2.6, Open results 2.7 and 2.8 and Theorems 2.9 and 2.10, Theorems 18.1 to 18.4 and Open results 18.5 and 18.6 and Theorem 35.1}{Theorem 41.1, Theorem 42.1, Theorem 44.1, Theorem 47.1 and Theorem 124.1}
The Born rule follows from the wash: the standing readout is the outcome probability by theorem, conditional on the axiom text and on two imports of its parent result, the Hilbert space of a qudit and the pair as a ray.
Every known quantity is a functional invariant under the matching group, the normalized trace over the register is the unique invariant state, so the canon, the trace of the projector against the state over the register dimension, is the unique readout, and the commutant of the matching group factors over the links.
An observer with reads at many sites on disjoint links sees commuting reads whose joint readout is the product of the one-site readouts, so the many-site observable is that product multiplied by the sign record; this rests on the monogamy lemma taken as a premise, on the valence bound taken as an input, and on the reading of the word sign as the phase of order $d$.
Distilling in $n$ copies projects onto a pair on the common support of dimension $d^n$, the product of the per-copy pairs, and the clock with the $d^n$-th roots and the shift with the carry satisfy the Weyl relation exactly and irreducibly at every pinned dimension and copy count, the lattice spacing falling by the factor $d$ at every further copy, with the ledger scale left open.
The observable exhaustion, whose one-site observer clause the many-site result lifts, is a further entry of this chapter, cited here for completeness.
\begin{chbridge}
The Born rule assigns to an outcome with projector $P$ on a state $\rho$ the probability $\mathrm{Tr}(\rho P)$, the only frame function of its kind on a Hilbert space of dimension at least three by the theorem of Gleason \cite{Gleason1957}, and Zurek derived it from envariance, the invariance of an entangled state under a unitary on the system that a unitary on its environment undoes: at a maximally entangled pair every unitary is undone, $(U\otimes\bar{U})\lvert\omega\rangle=\lvert\omega\rangle$, and these pairs of unitaries form the matching group of this chapter \cite{Zurek2003envariance,Zurek2005}; the evolution of a closed system is unitary, and a pure bipartite state has a maximally mixed reduced state on a support of rank $m$ iff its coefficient matrix is a multiple of a unitary on that support, the Schmidt lemma \cite{Nielsen2000}; the clock and shift operators satisfy the Weyl relation \cite{Weyl1927}. The correspondence is by name.
\end{chbridge}
\begin{chproof}[exact part and census part]
The premises of this derivation are the axiom and the six definitions of chapter~\ref{mod:FRAMEWORK}, the open result of chapter~\ref{mod:I-1-axiom-and-the-six-definitions}, the theorem of chapter~\ref{mod:I-3-records-reads-and-the-born-count-the-gain-lemma} and the theorem of chapter~\ref{mod:II-1-state-space-from-the-act-outputs}. The chain of premises runs back to the axiom through those chapters, and the longest such chain in this compendium is twelve chapters long.
The flat-reduced-state criterion is tested in both directions on the pinned family of pairs, the read transport, the join and the order step at every tabulated case; the commutant of the matching group is computed at a read and the unique invariant state found; reads at disjoint links and at one site are composed as operators; and the clock and the shift at dimension $d^n$ are built with the Weyl relation checked at every pinned pair; each theorem is conditional on the axiom text and on the open results it rests on.
\textit{The exact part.}
From the coefficient matrix $C$ of a pure pair on a common support of rank $m$, by the partial trace over the far end, we obtain the reduced state at the near end as the product $C C^{\dagger}$. The reduced state is flat on the support exactly when that product is a multiple of the support projector, which is the statement that $C$ is a multiple of a unitary on the support.
\begin{equation}
\rho_{\mathrm{near}} = C C^{\dagger} = \lambda\, P_{m}, \qquad \lambda = 1,\ \operatorname{rank} P_{m} = m .
\end{equation}
At the smallest instance, the pair ket of the two-qubit link, the support has rank $m = 2$, the scale $\lambda$ is one, the two support projectors have ranks two and two, the characteristic polynomial of $C C^{\dagger}$ has the flat form with two equal roots, and both directions of the equivalence hold, so the pair ket is a scaled partial isometry whose coefficient matrix is unitary on its support.
\begin{equation}
m = 2, \qquad \lambda = 1, \qquad \operatorname{rank}(P_{\mathrm{near}}, P_{\mathrm{far}}) = (2, 2), \qquad C C^{\dagger} = P_{2} .
\end{equation}
The unit central generator $G$ of the order step at the $(2,2)$ cell is the imaginary unit multiplied by the identity of the link ambient, the tensor product of the two ends, of dimension two multiplied by two, and its ambient trace is its trace over that ambient divided by the imaginary unit, the ambient dimension four, the per-end traces two and two. From $G$, by exponentiation at the twelve steps, we obtain the step as a unitary. The generator is anti-Hermitian and central at every one of its seven commutation tests, its Killing form vanishes, and its exponential at the twelve steps, the twelfth roots of unity multiplying the identity, is a unitary at every step with the group law met at all one hundred and forty-four products, the order of the exponential being twelve. A reader checks every line on the four-by-four identity by hand.
\begin{equation} \adjustbox{max width=\dimexpr\linewidth-3em\relax}{$\displaystyle \begin{gathered} G = i\,1_{2 \otimes 2}, \qquad G^{\dagger} = -G, \qquad \operatorname{tr}_{\mathrm{amb}}(G/i) = 2 \cdot 2 = 4, \qquad \kappa(G,G) = 0, \\ U_k = \exp\!\big(\tfrac{2\pi k}{12}\,G\big) = \zeta_{12}^{\,k}\,1_{2 \otimes 2},\ U_k U_k^{\dagger} = 1,\ U_{12} = 1 . \end{gathered}$} \end{equation}
From two pair kets sharing one qubit, by the contraction of their coefficient matrices at the shared end, we obtain the join at the triple of three qubits, the output is pure, it equals the pair ket on the composed support of rank two with scale one, the two factors are unitary on their supports, and the contraction is unitary on the composed support, so the join is a unitary.
\begin{equation}
\operatorname{rank}(C_{AB} C_{BC}) = 2, \qquad \lambda_{\mathrm{join}} = 1, \qquad (C_{AB} C_{BC})(C_{AB} C_{BC})^{\dagger} = P_{2} .
\end{equation}
\textit{The census part.}
By exhaustive enumeration of the 26870 matchings of the read transport (the standing monomial fourth-root groups at every end, generated by closure at the orders four, thirty-two, three hundred and eighty-four and six thousand one hundred and forty-four for the alphabet dimensions one to four, together with the whole unitary group at each end through its one hundred and sixty-six Lie generators, their Cayley unitaries and the Fourier matrix), the property tested at every matching being that the matching matrix is a scaled partial isometry with a scale-free Gram matrix and an orthogonal far set, every one of the 26870 matchings passes. The equivalence of the flat reduced state with the unitary coefficient matrix holds in both directions at every one of the sixteen pairs of the pinned family, the standing pair ket is a multiple of a unitary at every one of the four alphabet dimensions and every one of the six pair types with no differing type. The join at a shared qudit is a product of unitaries and a unitary on the composed support at every one of the sixty-four triples, the operator product of the pair-ket result re-met at every triple and the dressed joins canceling at the shared qudit at every one of the three hundred and twenty dressed cases. The order step exponentiates to a unitary at every one of the twelve steps at the link and at the cell, with the Weyl relation re-met at the conjugated letters at every step. The smallest case, the two-qubit matching, is the thirty-two monomial matchings of the first group above, every one a scaled partial isometry on the support of rank two, with two nonzero far-set vectors, and the four Lie generators anti-Hermitian with vanishing derivative of the Gram matrix.
The counter-tests alter the objects singly and read what moves. A non-flat pair planted among the pinned family moves the flatness and the unitarity cells at every one of the twelve planted readings. A non-unitary matching planted at the transport moves the partial-isometry, Gram and far-set cells. A join with no shared qudit is excluded by the join clause, as pinned. A degenerate triple reads coincident with the pair ket, as pinned, and the planted text alteration moves the four deciding cells. The matrices of the pair ket themselves are cells of the pair-ket result at its own chapter and are cited there, not re-printed here. The result was reproduced by a second engine built apart from the first, the two runs byte-identical at every one of the shared keys with no differing cell, and the five counter-tests and ten controls hold at their pinned readings.
The machine-checked package named at the reproducibility appendix certifies the statements of this chapter under the definitions of this compendium and proves no theorem of the literature.
\chcert{26870}
\end{chproof}
\begin{figure}[h]\centering
\setlength{\unitlength}{1pt}
\begin{picture}(260,150)
\qbezier(125.5,120.2)(107.1,88.2)(88.6,56.3)
\put(114.0,84.2){\makebox(0,0){\scriptsize $(2,2)$}}
\qbezier(93.1,48.5)(130.0,48.5)(166.9,48.5)
\put(130.0,56.5){\makebox(0,0){\scriptsize $(2,2)$}}
\put(130.0,128.0){\circle{14}}
\put(130.0,128.0){\makebox(0,0){\footnotesize $2$}}
\put(130.0,138.0){\makebox(0,0)[b]{\scriptsize $A$}}
\put(84.1,48.5){\circle{14}}
\put(84.1,48.5){\makebox(0,0){\footnotesize $2$}}
\put(84.1,58.5){\makebox(0,0)[b]{\scriptsize $B$}}
\put(175.9,48.5){\circle{14}}
\put(175.9,48.5){\makebox(0,0){\footnotesize $2$}}
\put(175.9,58.5){\makebox(0,0)[b]{\scriptsize $C$}}
\end{picture}
\caption{The smallest join of this chapter, two pair kets on two-qubit links sharing the qubit at the middle site, each site labeled by its dimension. The contraction at the shared site is the join, a unitary on the composed support of rank two.}
\end{figure}
\par\medskip\noindent{\itshape The rule and the step in hand, the next chapter draws the uncertainty relation from the relation the step obeys.}\label{narr:close:2-2}\par
\section{The uncertainty relation from the Weyl relation}\label{mod:II-3-uncertainty-relation-from-the-weyl-relation}
\chdesig{3-3}{2-3}{II-3-uncertainty-relation-from-the-weyl-relation}{II-3-uncertainty-relation-from-the-weyl-relation}
\label{narr:open:2-3}Theorem~41.1 obtains the uncertainty relation from the Weyl relation, resting on the state space and the step of the two preceding chapters.
\begin{chconditional}[The uncertainty relation from the Weyl relation]\label{box:II-3-uncertainty-relation-from-the-weyl-relation:1}
This result is conditional on the axiom text and on the results of the chapters named as its premises in the derivation below, each at its own standing.
Two reads of the one link in two bases do not commute: the commutator of the two label observables is a nonzero exact matrix at every pair of distinct bases at the dimensions two, three and four and zero at a basis against itself; the variance of every read at every register of the span equals the norm of the centred read; the product of the two variances is at least the squared modulus of half the commutator expectation at every pair of distinct bases and every register, transported by the whole standing group, with no violating pair and state at any dimension; and the saturating registers are the kernel of the Cauchy-Schwarz defect at every pair, the full register among them at none.
\end{chconditional}
\chstanding{a conditional theorem}
\chdeps{Definitions 1.1 to 1.6, Theorems 2.1 to 2.3, Open results 2.4 and 2.5, Theorem 2.6, Open results 2.7 and 2.8 and Theorems 2.9 and 2.10, Theorems 18.1 to 18.4 and Open results 18.5 and 18.6, Theorem 35.1 and Theorem 40.1}{no theorem beyond it at this release}
Two reads that commute only up to a sign form a further record whose order is the sign, a corollary of the axiom, the four definitions and the first read clause by the Weyl relation and named standard theorems with no step posed: the Weyl relation holds at every pair of letters at the dimensions two, three and four, the letter-shifted pair states form the Bell basis, the commutator phase at the qubit end is a sign and at the qutrit end a cube root, so no sign record stands at the qutrit end at any pair of reads; at the qubit end every anticommuting pair composes to one further operator with the order the sign, fixed under the standing group and every byproduct phase, and the three further records equal the exterior square of the traceless qubit space and the re-entry count of the two-point result.
\begin{chbridge}
For two observables $A$ and $B$ and a state $\rho$ the product of the two variances is at least the squared modulus of half the expectation of the commutator, $\Delta A\,\Delta B \geq \tfrac{1}{2}\,|\langle [A,B]\rangle|$, the Robertson form of the uncertainty relation, a corollary of the Cauchy-Schwarz inequality with equality iff the two centred operators applied to the state are proportional \cite{Robertson1929}. The correspondence is by name.
\end{chbridge}
\begin{chproof}[exact part and census part]
The premises of this derivation are the axiom and the six definitions of chapter~\ref{mod:FRAMEWORK}, the open result of chapter~\ref{mod:I-1-axiom-and-the-six-definitions}, the theorem of chapter~\ref{mod:I-3-records-reads-and-the-born-count-the-gain-lemma}, the theorem of chapter~\ref{mod:II-1-state-space-from-the-act-outputs} and the theorem of chapter~\ref{mod:II-2-born-rule-and-unitarity}. The chain of premises runs back to the axiom through those chapters, and the longest such chain in this compendium is twelve chapters long.
The two label observables are built at every tabulated pair of bases at the three dimensions and their commutator computed exactly, the Robertson bound is evaluated at every pair and register and transported by the standing group, and the saturating registers are found as the kernel of the Cauchy-Schwarz defect; the Weyl relation is computed at every pair of letters, the commutator phases are typed at both ends, and the further records are counted against the exterior square and the re-entry count. Both theorems are conditional on the axiom text, on the open results they rest on, and on results stated elsewhere.
\textit{The exact part.}
A basis of the one link at an end of dimension $d$ is a maximal commuting set of Weyl letters with its $d$ joint eigenprojectors of rank one, and the label observable of the basis is the operator carrying the label of each joint eigenprojector as its eigenvalue, the labels zero to $d$ less one in the labeling this result pins (the labels of a read a convention of the readout), so that at the qubit end a label observable is the rank-one projector onto the eigenvector labeled one, the Born face of this result and the one normalization at which every cell below is read. From the two label observables $A$ and $B$ of two distinct bases, by the matrix commutator, we obtain an exact nonzero matrix. At the smallest instance, the qubit end with the bases of the first and second Weyl letters, a reader writes the two projectors in the eigenbasis of the first letter, $A$ onto its eigenvector labeled one and $B$ onto the eigenvector of the second letter labeled one, and multiplies,
\begin{equation}
A = \begin{pmatrix} 0 & 0 \\ 0 & 1 \end{pmatrix}, \qquad B = \tfrac{1}{2}\begin{pmatrix} 1 & 1 \\ 1 & 1 \end{pmatrix}, \qquad [A, B] = \begin{pmatrix} 0 & -\tfrac{1}{2} \\ \tfrac{1}{2} & 0 \end{pmatrix} ,
\end{equation}
the entries of the commutator a half in modulus at the label normalization, and from the pair state at the full register of that end, by the Born readout (the trace of the projector against the flat state), we obtain the readout one half at every letter.
\begin{equation}
\langle A \rangle_{\mathrm{full}} = \langle B \rangle_{\mathrm{full}} = \tfrac{1}{2} .
\end{equation}
From the two variances of the reads at a register of the span, their product and the expectation of the commutator, by the Cauchy-Schwarz inequality applied to the centered reads, we obtain the bound with its defect at the same label normalization. At the register that is the first eigenstate of the first basis the variance of $A$ is zero, the variance of $B$ is a quarter (a projector with the expectation one half there, so its variance is a half less a quarter), the expectation of the commutator is zero since the commutator is off-diagonal, the squared covariance is zero and the gap is zero, so the bound saturates there,
\begin{equation}
\operatorname{Var}(A)\operatorname{Var}(B) \;\ge\; \Big|\tfrac{1}{2}\langle [A,B] \rangle\Big|^{2}, \qquad 0 \cdot \tfrac{1}{4} \;\ge\; 0 ,
\end{equation}
and the saturating registers at that pair of bases are the kernel of the Cauchy-Schwarz defect, six among the seven registers of the span, the full register not among them, at which both variances read a quarter, their product a sixteenth, and the gap a sixteenth.
\textit{The census part.}
By exhaustive enumeration of the pairs of distinct bases of the one link at the dimensions two, three and four (three, six and twenty-one unordered pairs, the bases the maximal commuting sets of Weyl letters with their joint eigenprojectors), the commutator of the two label observables is computed exactly at every pair and found nonzero at every one and zero at a basis against itself. The letter-face commutator is zero exactly at the Weyl phase one and nonzero at three, twenty-four and eighty-four letter pairs. Every bound is transported by every element of the standing group at each dimension (ninety-six, two thousand three hundred and four and eighty-four elements) with no violating pair and state. The variance of every read at every register of the span equals the norm of the centered read\, and the saturating registers are the kernel of the Cauchy-Schwarz defect at every pair, the full register among them at none. The smallest case is worked above at the qubit end.
The counter-tests alter the objects singly and read what moves. An altered pair of bases moves the Weyl and bound cells and the first posed step. An altered variance moves the variance and bound cells. The five counter-tests stand at their pinned letters and the ten controls read the same at both engines. The result was reproduced by a second engine built apart from the first, the two runs equal at every key.
The machine-checked package named at the reproducibility appendix certifies the statements of this chapter under the definitions of this compendium and proves no theorem of the literature.
\chcert{30}
\end{chproof}
\begin{figure}[h]\centering
\setlength{\unitlength}{1pt}
\begin{picture}(260,150)
\qbezier(86.0,75.0)(130.0,75.0)(174.0,75.0)
\put(130.0,83.0){\makebox(0,0){\scriptsize $(2,2)$}}
\put(77.0,75.0){\circle{14}}
\put(77.0,75.0){\makebox(0,0){\footnotesize $2$}}
\put(77.0,85.0){\makebox(0,0)[b]{\scriptsize $A$}}
\put(183.0,75.0){\circle{14}}
\put(183.0,75.0){\makebox(0,0){\footnotesize $2$}}
\put(183.0,85.0){\makebox(0,0)[b]{\scriptsize $B$}}
\end{picture}
\caption{The one link at its smallest instance, two qubit ends joined by the (2,2) link. The two reads of this chapter are two bases at one end, the labels of the Weyl letters.}
\end{figure}
\par\medskip\noindent{\itshape The thread turns from the relation between two reads to the act of one read, which is measurement.}\label{narr:close:2-3}\par
\section{Measurement: the wash at one act and the POVM}\label{mod:II-4-measurement-the-wash-at-one-act-and-the-povm}
\chdesig{3-4}{2-4}{II-4-measurement-the-wash-at-one-act-and-the-povm}{II-4-measurement-the-wash-at-one-act-and-the-povm}
\label{narr:open:2-4}Measurement in this construction is the wash at one act. Theorem~42.1 states the collective readout of the joined web of several bodies, the result the literature files under the positive operator-valued measure, resting on the records of chapter 1-3, the probe of chapter 1-6 and the step of chapter 2-2.
\begin{chconditional}[The collective readout of the joined web of several bodies]\label{box:II-4-measurement-the-wash-at-one-act-and-the-povm:1}
This result is conditional on the axiom text and on the results of the chapters named as its premises in the derivation below, each at its own standing.
At every pairing the docks admit, the collective readout of the joined web of $N$ bodies is defined at all six joined webs at $N$ two, three and four, the interference terms (the cross terms) between the copies exact and the coherence ratio an exact output at every one; at the reading named beside, the readout parts, the first parting the star at $N$ four refused at the valence bound.
\end{chconditional}
\chstanding{a conditional theorem}
\chdeps{Definitions 1.1 to 1.6, Theorems 18.1 to 18.4 and Open results 18.5 and 18.6, Theorems 28.1 and 28.2, Open results 28.3 and 28.4 and Theorem 28.5 and Theorem 40.1}{Theorem 24.1, Theorem 46.1, Theorem 47.1, Open result 166.5, Theorems 224.1 to 224.3 and Open result 224.4 and Theorem 229.1}
The standing web is one stabilizer state with two generators per link; distilling at every shared qudit is the entanglement swap, leaving the far ends a pair on their common support at every outcome of every one of the ten distinct triples (3412 swaps at 1100 shared qudits) with a Weyl letter as the byproduct; every tabulated read is a measurement in one basis at its slot, every sign record an adaptive dependency, the round count of the pattern equals the depth at every one of the 64 webs of the construction, and the output of the pattern is the normalized trace over the register at every one of the 309 reads.
Two open results stand here. The generalized measurement stands at its named clauses as a gloss: four of its five steps resolve at theorems (the wash at one act, the placement (the position of a tensor in a network), the rate as the washed square over the Born square, the dock register as the whole of the common support), but at the electron joined web the seventeen effects, each an exact nonnegative rational at the dock unit, sum to the cross class at the dock while the standing dock register is sixteen, so the completeness relation parts at the Born reading and at the count. The effects are built as operators on the dock register and their sum printed, but the identification of that sum with the identity of the register parts.
\begin{chbridge}
A generalized measurement is a set of positive operators $E_i$ summing to the identity with outcome probabilities $\mathrm{Tr}(\rho E_i)$, by the dilation theorem a projective measurement on a larger space \cite{Peres1993}; a measurement-based computation is a sequence of single-site measurements on a stabilizer state in rounds, the depth of the pattern the number of rounds \cite{Raussendorf2001}; and the entanglement swap projects two systems onto a maximally entangled pair and leaves their partners entangled up to a byproduct operator \cite{Zukowski1993}. The correspondence is by name: the compendium supplies a web of this theory read at its sinks, a doublet swap sign at the odd cycle, a Dirac form at an established generator and a finite record sum carrying the path-integral name at chapter~\ref{mod:II-1-state-space-from-the-act-outputs}, chapter~\ref{mod:II-5-indistinguishability-and-the-spin-statistics-sign}, chapter~\ref{mod:II-7-path-integral-as-the-finite-record-sum}, chapter~\ref{mod:III-5-spin-and-helicity}, the chapter titled The Dirac equation at the dock (chapter~\ref{mod:V-9-dirac-equation-at-the-dock}, among the instruments), chapter~\ref{mod:VI-2-w-and-the-z-and-their-masses}, chapter~\ref{mod:VI-2b-neutral-lepton-the-hub-unit-with-its-legs-open}, chapter~\ref{mod:VII-5-proton-to-electron-ratio-a-candidate}, chapter~\ref{mod:VIII-1-mass-unit-and-the-first-power}, chapter~\ref{mod:XIV-2-walls-census-in-physics-terms} and chapter~\ref{mod:XIV-7-what-remains-the-ungraded-page}, and the electron as the spin-one-half quantum of a field with its mass and spin as Lorentz labels corresponds to these objects by name.
\end{chbridge}
\begin{chproof}[exact part and census part]
The premises of this derivation are the axiom and the six definitions of chapter~\ref{mod:FRAMEWORK}, the theorem of chapter~\ref{mod:I-3-records-reads-and-the-born-count-the-gain-lemma}, the theorem of chapter~\ref{mod:I-6-read-at-a-register-the-probe-and-the-measurement} and the theorem of chapter~\ref{mod:II-2-born-rule-and-unitarity}. The chain of premises runs back to the axiom through those chapters, and the longest such chain in this compendium is twelve chapters long.
The joined webs of $N$ copies are enumerated at the pairings the two docks admit and the direct sum of the many-body result is computed at every one; the standing web is written as a stabilizer state, the swap is computed at every shared qudit and outcome, every read is typed as a basis measurement at its slot, and the round count is compared with the tabulated depth. The many-body theorem rests on the axiom text and the proved results it names alone; the pattern theorem is conditional on the open results it rests on.
\textit{The exact part.}
From the two-qubit pair type at the slot polarity, by the record spectrum, we obtain the record space of the pair, a roster of fifty-four, six swap-fixed records, twenty-four orbits, the four sectors holding twenty-four odd, eighteen mixed-even, twelve same-even and three closure records, fifty-seven records in all, the Born face ninety-six, the whole sixteen and the cut nine, with thirty one-exchange and three two-exchange terms. The two-exchange sum is an exact element of the ring,
\begin{equation}
G_{2}(2,2) = \tfrac{5}{216} + \tfrac{5}{216}\sqrt{2}, \qquad N = 24 + 18 + 12 + 3 = 57, \qquad Z = 96 .
\end{equation}
The smallest joined web of two bodies is two copies of the two-dock body (the body of the collective readout) paired at one dock. That body has nine sites, three qutrits $A$, $B$ and $C$, one site $H$ of dimension four, three qubits $a_{1}$, $a_{2}$ and $a_{3}$ and two dock qubits $w_{1}$ and $w_{2}$, and ten links, written by their two end sites, $AB$, $BC$ and $AC$ of type $(3,3)$, $HA$ of type $(3,4)$, $Ha_{1}$, $Ha_{2}$ and $Ha_{3}$ of type $(2,4)$, $a_{1}a_{2}$ of type $(2,2)$, and the two dock links $Bw_{1}$ and $Cw_{2}$ of type $(2,3)$. The measure of a link is the product of its two end dimensions, nine, twelve, eight, four and six at the five types. The two copies are joined by one pairing link $w_{1}w_{1}$ of type $(2,2)$ between their first docks, the unpaired dock links $Cw_{2}$ of both copies are struck, and the dressing of the joined web at two copies is nine. From these inputs, by the rule that the first power of a block is the dressing multiplied by the sum of the measures of its unstruck links, we obtain the first powers of the three blocks, each copy and the pairing link,
\begin{equation} \adjustbox{max width=\dimexpr\linewidth-3em\relax}{$\displaystyle \begin{gathered} S_{1}(c_{1}) = S_{1}(c_{2}) = 9 \cdot (9+9+9+12+8+8+8+4+6) = 657, \\ S_{1}(\text{pairing}) = 9 \cdot 4 = 36, \qquad S_{1} = 657 + 657 + 36 = 1350 , \end{gathered}$} \end{equation}
and from the first powers, by the rule that the pair count of two blocks is the product of their first powers taken in both orders, we obtain the four pair counts of the collective readout at the web face under the struck reading as exact integers, the pairs within one copy, the pairs at two links of different copies (the cross terms between the copies, carrying the sign minus), the pairs of a copy link with the pairing link, and the pairs of the pairing link with itself, their unsigned sum the square of the first power,
\begin{equation} \adjustbox{max width=\dimexpr\linewidth-3em\relax}{$\displaystyle \text{own} = 2 \cdot 657^{2} = 863298, \quad \text{cross} = -\,2 \cdot 657 \cdot 657 = -\,863298, \quad \text{cross-join} = 2 \cdot 2 \cdot 36 \cdot 657 = 94608, \quad \text{join-join} = 36^{2} = 1296, \quad 1350^{2} = 1822500 .$} \end{equation}
A reader reproduces the four counts with a few lines of arithmetic from the link list above.
\textit{The census part.}
By exhaustive enumeration of the pairings the two docks admit at two, three and four bodies, the six joined webs of the construction are built with every link typed at its slot (its passage, pair type and kind, two hundred and six links in all with the single body), the direct sum of the many-body readout is computed at every one, the cross terms between the copies are read as exact objects and the coherence ratio as an exact output at every web, the readout defined at all six. At the reading named beside, the readout parts, the first parting the star at four bodies refused at the valence bound. The smallest case, the two copies paired at their docks, is worked above.
The counter-tests alter the objects singly and read what moves. Striking the cross terms moves six cells, the first the double pair at two bodies. Leaving the docks unpaired moves the same six. The five counter-tests stand at their pinned letters, and the thirteen controls read the same at both engines at the spelling re-posed beside the pinned grammar (twelve at the pinned grammar, the one control re-posed at the sign census). The result was reproduced by a second engine built apart from the first, the two runs equal at every one of the cells of both.
The machine-checked package named at the reproducibility appendix certifies the statements of this chapter under the definitions of this compendium and proves no theorem of the literature.
\chcert{206}
\end{chproof}
\begin{figure}[h]\centering
\begin{adjustbox}{max width=\linewidth}\begin{tabular}{llrrrr}\hline bodies & web & own pairs & cross pairs & cross-join & join-join \\ \hline 2 & the double pair & 863298 & 863298 & 94608 & 1296 \\ \hline \end{tabular}\end{adjustbox}
\caption{The smallest joined web of this chapter as a listed table, two copies of the body paired at their docks, with the four pair counts of its collective readout at the web face under the struck reading, every count an exact integer.}
\end{figure}
\par\medskip\noindent{\itshape With measurement read as a wash, the next chapter asks what the construction says about bodies that cannot be told apart.}\label{narr:close:2-4}\par
\section{Indistinguishability and the spin-statistics sign}\label{mod:II-5-indistinguishability-and-the-spin-statistics-sign}
\chdesig{3-5}{2-5}{II-5-indistinguishability-and-the-spin-statistics-sign}{II-5-indistinguishability-and-the-spin-statistics-sign}
\label{narr:open:2-5}Theorem~43.1 obtains indistinguishability from the twirl, with the sign of the many-body exchange read at the same place, resting on the link, the moves, the probe and the state space.
\begin{chconditional}[Indistinguishability from the twirl]\label{box:II-5-indistinguishability-and-the-spin-statistics-sign:1}
This result is conditional on the axiom text and on the results of the chapters named as its premises in the derivation below, each at its own standing.
Under the reading that takes the two copies of a site as one qudit of dimension $d^2$, the standing two-copy configuration is built at all eighteen objects, the swap is washed by the twirl of the identification group, and the readout is symmetric under the swap; the sign under the swap is plus one at the amplitude level at every object, against the odd-cycle sign minus one of the exchange triangle, which holds beside and obstructs nothing.
\end{chconditional}
\chstanding{a conditional theorem}
\chdeps{Definitions 1.1 to 1.6, Theorems 9.1 to 9.7 and Open result 9.8, Theorems 28.1 and 28.2, Open results 28.3 and 28.4 and Theorem 28.5, Theorem 35.1 and Open results 241.1 and 241.2, Theorems 241.3 and 241.4 and Open results 241.5 to 241.7}{Theorem 46.1 and Theorem 47.1}
The exchange of the two-dimensional objects of a body (the doublet swap) at a generic qubit read is realised by transpositions at each of the eighteen objects, and its sign is minus one to the doublet count, nine objects at plus one and nine at minus one, in agreement with the Standard Model statistics names; this is a conditional theorem, the second fallback of the open statistics results.
Two pairs sharing a direction are exchanged by a basis identification that fixes nothing, so the shared direction is washed and carries no read: the exclusion of the $d$ plus first pair is a corollary of the counting at every one of the 1691 links of the construction under both readings of the basis, and no-cloning is the same corollary at one pair, the copy excluded by the count at 1229 among the 1334 sites and by the wash at the remaining 105.
Two open results stand here. The standing exchange cycle is odd at every object while the doublet count is one at the nine chain bodies alone, so the clause odd cycle iff one doublet parts at nine of the eighteen objects, the W word first. The doublet swap carries the sign minus one to the doublet count at every object and is an admissible identification at the own readable reads of fifteen objects, the gluon word first among the three that part.
\begin{chbridge}
The state of $N$ identical particles is symmetric or antisymmetric under every exchange of two of them, and the spin-statistics theorem ties the sign to the spin \cite{Pauli1940}; the exclusion principle is the antisymmetric case at one orbital, and no unitary copies an unknown state onto a second system \cite{WoottersZurek1982}; the exchange of two subsystems of dimension $d$ is a permutation of the $d^2$ product labels with $d(d-1)/2$ transpositions \cite{Nielsen2000}. The correspondence is by name.
\end{chbridge}
\begin{chproof}[exact part and census part]
The premises of this derivation are the axiom and the six definitions of chapter~\ref{mod:FRAMEWORK}, the theorem of chapter~\ref{mod:I-2-link-and-its-six-genres}, the theorem of chapter~\ref{mod:I-4-moves-in-one-table}, the theorem of chapter~\ref{mod:I-6-read-at-a-register-the-probe-and-the-measurement} and the theorem of chapter~\ref{mod:II-1-state-space-from-the-act-outputs}. The chain of premises runs back to the axiom through those chapters, and the longest such chain in this compendium is twelve chapters long.
The two copies of every object are built at the doubled site, the swap taken as an identification there, the twirl applied and the readout compared under the swap; the generic swap is built at a generic qubit read and its sign computed at the abstract, label and determinant faces; and the exchange of two pairs sharing a direction is built as an identification, its fixed directions counted and the exclusion compared with the counting at every link under both readings. The theorems are conditional on the open results they rest on, among them the Born rule from the wash and the holonomy sign definition, and on the roster of the construction.
\textit{The exact part.}
From the two copies of a site of dimension $d$ read as one qudit of dimension $d^{2}$, by the doubled pair at the two-qubit type, we obtain the support two, the doubled support four and the two-copy register of dimension sixteen. The swap of the copies fixes the two-copy pair ket with the amplitude sign plus one and the permutation sign minus one, at Schmidt rank four.
\begin{equation}
r = 2, \qquad D = 4, \qquad \dim R = 16, \qquad \sigma_{\mathrm{amp}} = +1, \qquad \sigma_{\mathrm{perm}} = -1 .
\end{equation}
The fixed set at a type is the commutant of the identification group on the two-copy register, the operators the twirl leaves fixed, spanned at the doubled support by the orbits of the adjacent transpositions under the monomial subgroup (the commutant of that subgroup of dimension three, the commutant of the whole group of dimension two inside it), with the basis the identity, the pair projector $\Pi$ (the projector of the two-copy pair ket on the support) and the three operators $O_{1}$, $O_{2}$ and $O_{3}$ as the twirl tables of this result list them, and the readout of an operator $M$ at the configuration is the trace of $M$ multiplied by $\Pi$ over the register dimension, sixteen at this type. From the five readouts of the twirled configuration at that type, by that trace, we obtain the exact rationals one sixteenth, one sixteenth, one sixteenth, three sixteenths and zero, symmetric under the swap at every one of the five, the first two a reader checks by hand (the trace of a rank-one projector is one and the pair projector squares to itself) and the fourth the doubled support less one over the register,
\begin{equation}
\langle M \rangle = \frac{\operatorname{tr}(M\,\Pi)}{16}, \qquad \langle 1 \rangle = \tfrac{1}{16}, \quad \langle \Pi \rangle = \tfrac{1}{16}, \quad \langle O_{1} \rangle = \tfrac{1}{16}, \quad \langle O_{2} \rangle = \tfrac{D - 1}{16} = \tfrac{3}{16}, \quad \langle O_{3} \rangle = 0 .
\end{equation}
From the exchange cycle of the standing configuration, the triangle of length three, by the parity of its length, we obtain the odd-cycle sign minus one, which stands beside the swap sign plus one and obstructs nothing.
\begin{equation}
\sigma_{\mathrm{cycle}} = (-1)^{3} = -1, \qquad \sigma_{\mathrm{swap}} = +1 .
\end{equation}
\textit{The census part.}
By exhaustive enumeration of the eighteen objects of the construction (the thirteen bodies and the five carriers, re-met in order against their tables), the two-copy configuration of each is built at the doubled sites and doubled pairs, the swap taken as a basis identification there, washed by the twirl of the identification group, and the readout compared under the swap, the swap sign is plus one at the amplitude level at every one of the eighteen, the readouts are symmetric at every pair, and the odd-cycle sign is minus one at every object. The smallest case is the carrier with five sites and five links, doubled to ten sites and ten pairs, its hub of dimension four, its qutrit and its three qubits, the swap fixes all five pairs, the label twirl vanishes at all five and the readouts are symmetric at all five, the swap sign plus one.
The counter-tests alter the objects singly and read what moves. The first moves three hundred and forty-six cells among them the wash and symmetry cells at the two-qubit type. The second moves five hundred and six cells among them the copy and symmetry cells and the first failing object. The five counter-tests stand at their pinned letters and the eleven controls read the same at both engines and equal their expectations. The result was reproduced by a second engine built apart from the first, the two runs equal at every key.
The machine-checked package named at the reproducibility appendix certifies the statements of this chapter under the definitions of this compendium and proves no theorem of the literature.
\chcert{18}
\end{chproof}
\begin{figure}[h]\centering
\setlength{\unitlength}{1pt}
\begin{picture}(260,150)
\qbezier(122.7,122.7)(104.8,109.7)(86.9,96.7)
\put(109.5,103.2){\makebox(0,0){\scriptsize $(3,4)$}}
\qbezier(82.4,82.8)(89.2,61.8)(96.1,40.7)
\put(96.8,64.2){\makebox(0,0){\scriptsize $(2,4)$}}
\qbezier(86.9,86.1)(120.4,61.8)(153.9,37.4)
\put(125.1,68.2){\makebox(0,0){\scriptsize $(2,4)$}}
\qbezier(88.6,91.4)(130.0,91.4)(171.4,91.4)
\put(130.0,99.4){\makebox(0,0){\scriptsize $(2,4)$}}
\qbezier(107.8,32.1)(130.0,32.1)(152.2,32.1)
\put(130.0,40.1){\makebox(0,0){\scriptsize $(2,2)$}}
\put(130.0,128.0){\circle{14}}
\put(130.0,128.0){\makebox(0,0){\footnotesize $3$}}
\put(130.0,138.0){\makebox(0,0)[b]{\scriptsize $Q$}}
\put(79.6,91.4){\circle{14}}
\put(79.6,91.4){\makebox(0,0){\footnotesize $4$}}
\put(79.6,101.4){\makebox(0,0)[b]{\scriptsize $H$}}
\put(98.8,32.1){\circle{14}}
\put(98.8,32.1){\makebox(0,0){\footnotesize $2$}}
\put(98.8,42.1){\makebox(0,0)[b]{\scriptsize $q_1$}}
\put(161.2,32.1){\circle{14}}
\put(161.2,32.1){\makebox(0,0){\footnotesize $2$}}
\put(161.2,42.1){\makebox(0,0)[b]{\scriptsize $q_2$}}
\put(180.4,91.4){\circle{14}}
\put(180.4,91.4){\makebox(0,0){\footnotesize $2$}}
\put(180.4,101.4){\makebox(0,0)[b]{\scriptsize $q_3$}}
\end{picture}
\caption{The smallest object of this chapter, a carrier web of five sites and five links, a qutrit, a hub of dimension four and three qubits, each site labeled by its dimension and each link by its pair type. Its two-copy configuration doubles every site and every pair.}
\end{figure}
\par\medskip\noindent{\itshape From the exchange of bodies the thread passes to the quantity conserved along the step, energy.}\label{narr:close:2-5}\par
\section{Energy and the Hamiltonian: the first-power operator and the step}\label{mod:II-6-energy-and-the-hamiltonian-the-first-power-operator-and-the-step}
\chdesig{3-6}{2-6}{II-6-energy-and-the-hamiltonian-the-first-power-operator-and-the-step}{II-6-energy-and-the-hamiltonian-the-first-power-operator-and-the-step}
\label{narr:open:2-6}Theorem~44.1 reads energy as the first-power operator and the step, the object the literature calls the Hamiltonian appearing here as the generator of the step read at the first power.
\begin{chconditional}[Energy as the first-power operator and the step]\label{box:II-6-energy-and-the-hamiltonian-the-first-power-operator-and-the-step:1}
This result is conditional on the axiom text and on the results of the chapters named as its premises in the derivation below, each at its own standing.
The first-power chain sum is an operator on the record space of every one of the 309 tabulated reads with the first-power sum of the web as its eigenvalue; it commutes with the step, whose generator (the step generator) is re-met at every pair type with one level and conserved at all 32851 census moves; and the product generator carries the first-power sum as its level, the Born square of the body level equal to the standing rest energy at all thirteen bodies, so the Schrodinger form holds at the record as a composition of two established objects.
\end{chconditional}
\chstanding{a conditional theorem}
\chdeps{Definitions 1.1 to 1.6, Theorems 18.1 to 18.4 and Open results 18.5 and 18.6, Theorem 40.1, Open results 241.1 and 241.2, Theorems 241.3 and 241.4 and Open results 241.5 to 241.7 and Open results 245.1 and 245.2, Theorem 245.3, Open results 245.4 to 245.8 and Theorem 245.9}{Open result 45.1, Theorem 47.1, Theorem 146.1, Theorem 225.1 and Theorem 226.1 and Open result 226.2}
The energy of a body is its chain sum at the first power, exact at the nine bodies and the 255 tabulated joins and additive in rank with the dressing a common factor at the 43 non-interacting joins, and the standing mass office is its Born square, so energy and mass are identified through the square at the register.
At every one of the 69211 registers of the 2276 tabulated webs and sinks the step generator equals the order operator in the standing unit, exactly, which dissolves the order-operator question.
The arrow transfer operator of every web of the construction satisfies the reversal identity, its 4080 prime cycles to the cap carrying a discrete real-face argument, zero at every even cycle and a half turn at the odd cycles of two-basis support. At a star through a centre of degree three the 169 records, their split 79, 48 and 42 and the closure 3 are re-met exactly under the record rule, a conditional theorem whose other face is open. The site-determinant identity holds for the register-weighted arrow operator at all 64 webs of the construction.
Three open results stand here. The standing generator, a global phase with one level, equals the first-power sum at the unit at none of the 309 reads. The standing transport parts from the exponential of the first-power operator at one tick at all 1692 links, the first-power operator generating a phase alone; of five corollary classes beside, two are derived and three are glosses. The exponent spectrum of the transport generator parts from the standing order spectrum at all 1692 links, a parting of dimension dissolved at the second result of its line.
\begin{chbridge}
The evolution of a closed system is a one-parameter group of unitaries, by the theorem of Stone the exponential of a self-adjoint generator, the Hamiltonian, whose eigenvalues are the energy levels \cite{Stone1932}; the rest energy of a body is its mass at the unit where the speed of light is one \cite{Dirac1958}; and the determinant of one minus $uB$ for the non-backtracking operator $B$ of a finite graph collapses to a determinant over the sites \cite{Bass1992}. The correspondence is by name: the compendium supplies an isometry algebra of the Hermitian split at the Lorentz brackets, a lattice point group, boosts from the order register and a transverse traceless part at the open read, every one a finite or algebraic object of this theory carrying the group name at chapter~\ref{mod:III-2-rotation-group}, chapter~\ref{mod:III-3-lorentz-group-and-the-signature}, chapter~\ref{mod:III-4-boosts-and-the-one-unit}, chapter~\ref{mod:III-5-spin-and-helicity}, the chapter titled Causal structure and general covariance (chapter~\ref{mod:III-8-causal-structure-and-general-covariance}, among the instruments) and chapter~\ref{mod:XIV-7-what-remains-the-ungraded-page}, a web of this theory read at its sinks, a doublet swap sign at the odd cycle, a Dirac form at an established generator and a finite record sum carrying the path-integral name at chapter~\ref{mod:II-1-state-space-from-the-act-outputs}, chapter~\ref{mod:II-5-indistinguishability-and-the-spin-statistics-sign}, chapter~\ref{mod:II-7-path-integral-as-the-finite-record-sum}, chapter~\ref{mod:III-5-spin-and-helicity}, the chapter titled The Dirac equation at the dock (chapter~\ref{mod:V-9-dirac-equation-at-the-dock}, among the instruments), chapter~\ref{mod:VI-2-w-and-the-z-and-their-masses}, chapter~\ref{mod:VI-2b-neutral-lepton-the-hub-unit-with-its-legs-open}, chapter~\ref{mod:VII-5-proton-to-electron-ratio-a-candidate}, chapter~\ref{mod:VIII-1-mass-unit-and-the-first-power}, chapter~\ref{mod:XIV-2-walls-census-in-physics-terms} and chapter~\ref{mod:XIV-7-what-remains-the-ungraded-page} and a finite re-matching group at the known clause, a representation census at the open ends and the sinks, an attach-site basis carrying the generation name and a sink class carrying the doublet name at chapter~\ref{mod:IV-1-gauge-invariance-as-the-known-clause}, chapter~\ref{mod:IV-3-representations-from-the-open-end-census}, chapter~\ref{mod:IV-4-cartan-weights-and-the-charge}, chapter~\ref{mod:IV-6-conservation-as-the-kernel-of-the-move-matrix}, chapter~\ref{mod:IV-7-noether-at-the-sinks-and-the-kernel-identification}, chapter~\ref{mod:IV-9-one-rule-of-a-read-and-the-word-table}, chapter~\ref{mod:VI-3-doublets-handedness-and-parity-violation}, chapter~\ref{mod:VII-7-quark-content-and-flavour}, chapter~\ref{mod:VIII-4-three-generations-as-attach-sites}, chapter~\ref{mod:XIV-7-what-remains-the-ungraded-page}, and the mass named as the Casimir label of a Lorentz representation and the rest eigenvalue of a Hamiltonian, the field named as a spinor field with the mass term the literature assigns it and a generation as a copy of the representation content of the symmetry correspond to these objects by name.
\end{chbridge}
\begin{chproof}[exact part and census part]
The premises of this derivation are the axiom and the six definitions of chapter~\ref{mod:FRAMEWORK}, the theorem of chapter~\ref{mod:I-3-records-reads-and-the-born-count-the-gain-lemma}, the theorem of chapter~\ref{mod:I-4-moves-in-one-table}, the theorem of chapter~\ref{mod:I-5-saturation-in-one-figure-seeds-waves-termination-sinks} and the theorem of chapter~\ref{mod:II-2-born-rule-and-unitarity}. The chain of premises runs back to the axiom through those chapters, and the longest such chain in this compendium is twelve chapters long.
The first-power operator is built from the record-level gain operator and its eigenvalue, commutator and product generator computed at every read; the chain sum of every body and join is computed at the first power and compared with the mass office; the step generator and the order operator are compared at every register; and the arrow operators are built at the register and letter faces with their identities, cycles and determinants computed, each theorem conditional on the axiom text and on the open results it rests on.
\textit{The exact part.}
From the measure of a web and its hub product, by the first-power rule with nothing struck, we obtain the first-power sum of the web. At the smallest read of the catalog, the read of the web of measure one hundred and sixteen at its two qubit sites, the hub product is thirty-six and the first-power sum is
\begin{equation}
S_{0} = 116 \cdot 36 = 4176 .
\end{equation}
From the record space of that read at the slot reading, by the record-level gain operator, we obtain fifty-seven records with the gains one at twenty-four, two at thirty and four at three, the trace of the gain operator an exact element of the ring equal to its closed form, and the Born face ninety-six.
\begin{equation}
\operatorname{tr}\, G = 30 + 30\sqrt{2}, \qquad Z = 96, \qquad N = 24 + 30 + 3 = 57 .
\end{equation}
From the first-power operator at that read, by its eigenvalue on the record space, we obtain the first-power sum of the web as the eigenvalue with multiplicity fifty-seven, its commutator with the step generator zero at every one of the twelve commutation tests, and its level at the hub unit and at the dock unit the exact rationals
\begin{equation}
H\,\psi = 4176\,\psi, \qquad [H, G_{\mathrm{step}}] = 0, \qquad \ell_{\mathrm{hub}} = \tfrac{1392}{67}, \qquad \ell_{\mathrm{dock}} = \tfrac{4176}{67} .
\end{equation}
From the electron body with its two struck links, by the body rule, we obtain the dock-struck rank sixty-seven, the hub dressing three and the body level two hundred and one, whose Born square at the hub unit is one and at the dock unit nine, equal to the standing rest energy at both units.
\begin{equation}
S = 67 \cdot 3 = 201, \qquad \ell_{\mathrm{hub}} = 1,\ \ell_{\mathrm{dock}} = 3, \qquad \ell_{\mathrm{hub}}^{2} = 1,\ \ell_{\mathrm{dock}}^{2} = 9 .
\end{equation}
\textit{The census part.}
By exhaustive enumeration of the 309 tabulated reads of the catalog (every read of a web of the construction at its record space), the first-power operator is built at each from the record-level gain operator and its eigenvalue found equal to the first-power sum of the web at every read. The step generator is re-met at every one of the six pair types with one level, and its central functional is conserved at every one of the 32851 census moves. The commutator of the first-power operator with the step generator is zero at every read, and the product generator carries the first-power sum as its level at every read, with the Born square of the body level equal to the standing rest energy at every one of the thirteen bodies at both units. The smallest case, the read of the web of measure one hundred and sixteen, is worked above.
The counter-tests alter the objects singly and read what moves, every one at its pinned letter at both engines, and the controls (the webs, the bodies, the reads, the sign result, the flat-pair result and the rest-energy result re-met by hash) read the same at both engines. The result was reproduced by a second engine built apart from the first, the two runs equal at every key.
The machine-checked package named at the reproducibility appendix certifies the statements of this chapter under the definitions of this compendium and proves no theorem of the literature.
\chcert{32851}
\end{chproof}
\begin{figure}[h]\centering
\setlength{\unitlength}{1pt}
\begin{picture}(260,150)
\qbezier(130.0,119.0)(130.0,75.0)(130.0,31.0)
\put(138.0,75.0){\makebox(0,0){\scriptsize $(3,3)$}}
\qbezier(123.6,121.6)(103.5,101.5)(83.4,81.4)
\put(109.2,95.8){\makebox(0,0){\scriptsize $(2,3)$}}
\qbezier(136.4,28.4)(156.5,48.5)(176.6,68.6)
\put(150.8,54.2){\makebox(0,0){\scriptsize $(2,3)$}}
\put(130.0,128.0){\circle{14}}
\put(130.0,128.0){\makebox(0,0){\footnotesize $3$}}
\put(130.0,138.0){\makebox(0,0)[b]{\scriptsize $B$}}
\put(77.0,75.0){\circle{14}}
\put(77.0,75.0){\makebox(0,0){\footnotesize $2$}}
\put(77.0,85.0){\makebox(0,0)[b]{\scriptsize $w_1$}}
\put(130.0,22.0){\circle{14}}
\put(130.0,22.0){\makebox(0,0){\footnotesize $3$}}
\put(130.0,32.0){\makebox(0,0)[b]{\scriptsize $C$}}
\put(183.0,75.0){\circle{14}}
\put(183.0,75.0){\makebox(0,0){\footnotesize $2$}}
\put(183.0,85.0){\makebox(0,0)[b]{\scriptsize $w_2$}}
\end{picture}
\caption{The electron body of this chapter at its two docks, the core link between the two qutrit sites and the two dock links to the qubit sites, each site labeled by its dimension. The two dock links are the struck links of the body rule, the dock-struck rank sixty-seven and the hub dressing three giving the body level two hundred and one.}
\end{figure}
\par\medskip\noindent{\itshape With the step and its generator in hand, the next chapter sums over the records the step leaves.}\label{narr:close:2-6}\par
\section{The path integral as the finite record sum}\label{mod:II-7-path-integral-as-the-finite-record-sum}
\chdesig{3-7}{2-7}{II-7-path-integral-as-the-finite-record-sum}{II-7-path-integral-as-the-finite-record-sum}
\label{narr:open:2-7}Open result~45.1 states, as an open result with its parting named, the sum over paths as the finite sum over records, resting on the records, the moves, the saturation and the energy chapter.
\begin{chwall}[The path integral as the finite record sum]\label{box:II-7-path-integral-as-the-finite-record-sum:1}
The path integral as the finite record sum is stated at its five components and re-met at the established tables, the grading at the orders zero to three, the measure at 140 links of sixteen webs, the termination at order three, and the record space of the one link at 172 letters with the support rank 319 equal to the Born count; yet the first entry of its keying, the chiral anomaly at the clause of the anomaly sums over the sixteen words, is unkeyed, the eighteen entries beside keyed.
Open are the five results above, the path integral as the finite record sum with its first entry unkeyed the open object of this chapter; the read order, the read at every order and the read through the body stand as theorems. Re-posed at the first unkeyed clause: the chiral anomaly entry stays a gloss with its cause named, the sixteen derived rows of the open-end census absent at the reading of every recorded census, the second census at eight derived rows and the first at fourteen with the gate of the census unmet at two of sixteen at both, so that the anomaly sums over the sixteen words stay open at this chapter; the sums over the derived words as they stand are printed beside, the cubic sum minus 675, the mixed sum at the doublet rows minus 3, the mixed sum at the colour rows 0 and the linear sum minus 21 at the second census and 402, minus 6, 24 and 24 at the first, every one an exact integer of the sixths lattice compared with nothing.
\end{chwall}
\chstanding{an open result}
\chdeps{Definitions 1.1 to 1.6, Theorems 18.1 to 18.4 and Open results 18.5 and 18.6, Theorem 44.1, Open results 241.1 and 241.2, Theorems 241.3 and 241.4 and Open results 241.5 to 241.7 and Open results 245.1 and 245.2, Theorem 245.3, Open results 245.4 to 245.8 and Theorem 245.9}{Theorem 146.1}
The order of a read at the (2,3) link is the number of signed exchanges its record survives, zero at the 79 odd records, one at the 90 one-exchange records and two at the 3 closure records, equal to the number of factors of the gain count at every one of the 172.
The read at every order is a finite series terminating at order two, the order-one read at the one link equal to the first-order record of the gain result at all 172 records, the order-two read equal to the second-order cells of the held link, and at the held pair the order-$k$ term at rank $n$ the term at rank one over $n$ squared at all eight ranks; a conditional theorem.
Through a body the read at every order is a series of two exact terms at every one of the thirteen bodies of the construction, the order-one read the first-power chain sum and the order-two read an exact element of the field of root two and root three, nonzero at eight of the nine chain bodies, the orders beyond two zero; a conditional theorem.
Four further open results stand here. The Born square of the body chain sum over the electron chain sum meets an established monomial face at none of the eight bodies beyond the electron. The body read of the charter equals the standing rest-energy invariant at none of the nine bodies, the mass office differing at the eight non-unit bodies by one missing factor, one over three squared. The path count at the record acts parts from the standing correspondence at the first root and the first wave. At the pinned sample of 3438 seeded sinks the forward paths of the move graph are finitely many at every sink, but the path sum equals the first power at two sinks alone.
\begin{chbridge}
The amplitude between two configurations is a sum over paths of the exponential of the action, for a finite system a finite sum of products of transfer amplitudes \cite{Feynman1948}; a theory whose perturbative expansion terminates at a finite order carries finitely many terms \cite{Peskin1995}; and the Rydberg ladder falls as the inverse square of the rank \cite{Dirac1958}. The correspondence is by name: the compendium supplies a joined saturation carrying the hydrogen name, a two-vertex term carrying the Rydberg name, cells at the hydrogen sink carrying the Lamb and hyperfine names, self-terms carrying the coupling names and joined electron sinks carrying the periodic-table name at chapter~\ref{mod:IV-8-ward-identity-and-the-anomaly-sums}, chapter~\ref{mod:V-2-electron-anomalous-dipole-factor-through-order-three}, the chapter titled The muon anomalous dipole factor (chapter~\ref{mod:V-3-muon-anomalous-dipole-factor}, among the instruments), chapter~\ref{mod:V-4-hydrogen-the-rydberg-ladder-and-the-fine-structure}, chapter~\ref{mod:V-5-lamb-shift-and-the-casimir-effect}, chapter~\ref{mod:V-8-photon-docking-the-join-face-and-pair-production}, the chapter titled The mesons and the nucleon at the cap (chapter~\ref{mod:VII-9-mesons-and-the-nucleon-at-the-cap}, among the instruments), chapter~\ref{mod:XII-3-chain-of-pairs-and-the-charged-periodic-sink}, chapter~\ref{mod:XII-5-periodic-table-as-the-joined-electron-sinks} and chapter~\ref{mod:XIV-7-what-remains-the-ungraded-page}, a step generator at the order coordinate, a first-power operator, a direction functional and an office at the common support carrying the energy, momentum and mass names at chapter~\ref{mod:I-3-records-reads-and-the-born-count-the-gain-lemma}, chapter~\ref{mod:I-6-read-at-a-register-the-probe-and-the-measurement}, chapter~\ref{mod:II-6-energy-and-the-hamiltonian-the-first-power-operator-and-the-step}, chapter~\ref{mod:III-1-order-coordinate-and-its-theorem}, chapter~\ref{mod:III-6-p-c-t-and-cpt}, chapter~\ref{mod:IV-6-conservation-as-the-kernel-of-the-move-matrix}, chapter~\ref{mod:V-7-running-of-the-coupling}, the chapter titled The Dirac equation at the dock (chapter~\ref{mod:V-9-dirac-equation-at-the-dock}, among the instruments), chapter~\ref{mod:VIII-1-mass-unit-and-the-first-power}, chapter~\ref{mod:VIII-5-unified-mass-readout}, chapter~\ref{mod:X-2-field-equation-at-the-counting-face}, the chapter titled The gravitational and cosmological constants table (chapter~\ref{mod:X-9-gravitational-and-cosmological-constants-table}, among the instruments), chapter~\ref{mod:XI-1-second-law-as-the-boundary-count-under-acts} and chapter~\ref{mod:XIV-7-what-remains-the-ungraded-page} and an isometry algebra of the Hermitian split at the Lorentz brackets, a lattice point group, boosts from the order register and a transverse traceless part at the open read, every one a finite or algebraic object of this theory carrying the group name at chapter~\ref{mod:III-2-rotation-group}, chapter~\ref{mod:III-3-lorentz-group-and-the-signature}, chapter~\ref{mod:III-4-boosts-and-the-one-unit}, chapter~\ref{mod:III-5-spin-and-helicity}, the chapter titled Causal structure and general covariance (chapter~\ref{mod:III-8-causal-structure-and-general-covariance}, among the instruments) and chapter~\ref{mod:XIV-7-what-remains-the-ungraded-page}, and the constant named as the scale of the Coulomb spectrum of hydrogen, the energy named as the eigenvalue of a Hamiltonian generating a continuous evolution and the mass named as the Casimir label of a Lorentz representation and the rest eigenvalue of a Hamiltonian correspond to these objects by name.
\end{chbridge}
\begin{chproof}[census]
The premises of this derivation are the axiom and the six definitions of chapter~\ref{mod:FRAMEWORK}, the theorem of chapter~\ref{mod:I-3-records-reads-and-the-born-count-the-gain-lemma}, the theorem of chapter~\ref{mod:I-4-moves-in-one-table}, the theorem of chapter~\ref{mod:I-5-saturation-in-one-figure-seeds-waves-termination-sinks} and the theorem of chapter~\ref{mod:II-6-energy-and-the-hamiltonian-the-first-power-operator-and-the-step}. The chain of premises runs back to the axiom through those chapters, and the longest such chain in this compendium is twelve chapters long.
The order of every record of the (2,3) link is read as the count of signed exchanges it survives and compared with the factors of the gain count; the order-one and order-two reads at the one link are compared with the gain result and the held link; and the order-two read of every body is computed through the dark ports and the hubs at the cut register; each theorem is conditional on the open results it rests on, the two order theorems besides on conditional theorems at their own statements.
By exhaustive enumeration of the nineteen entries of the keying (the quantum-field-theory terms of the literature that the finite record sum is asked to key), the object is stated once at its five components and re-met at the established tables, the grading at the orders zero to three over eighteen webs, the measure at one hundred and forty links of sixteen webs, the termination at order three, and the record space of the one link at one hundred and seventy-two letters with the support rank three hundred and nineteen equal to the Born count. The property tested at every entry is that it is keyed at a clause of the object by hash with its fragment found and its checks passing. Eighteen entries are keyed and one is not, the first unkeyed entry the chiral anomaly at the clause of the anomaly sums over the sixteen words, which names the parting of this open result. The smallest case is the first entry, the Fock space as the many-particle state space, keyed at the first clause (the many-particle space as the direct sum over records by order) at the grading component, its fragment found and all eighteen of its checks passing.
The counter-tests alter the objects singly and read what moves. Striking the grading moves twenty-six cells, the first the grading orders of the object. The three counter-tests stand at their pinned letters and the ten controls read the same at both engines. A second engine built apart from the first reproduced every cell.
The machine-checked package named at the reproducibility appendix certifies the statements of this chapter under the definitions of this compendium and proves no theorem of the literature.
\chcert{19}
\end{chproof}
\par\medskip\noindent{\itshape The sum over records leads to the phase-space picture of the next chapter.}\label{narr:close:2-7}\par
\section{The Wigner function and the many-body sign}\label{mod:II-8-wigner-function-and-the-many-body-sign}
\chdesig{3-8}{2-8}{II-8-wigner-function-and-the-many-body-sign}{II-8-wigner-function-and-the-many-body-sign}
\label{narr:open:2-8}Theorem~46.1 reads the phase-space function of the literature and the many-body sign together, resting on the records, the measurement of chapter 2-4 and the exchange of chapter 2-5.
\begin{chconditional}[The Wigner function and the many-body sign]\label{box:II-8-wigner-function-and-the-many-body-sign:1}
This result is conditional on the axiom text and on the results of the chapters named as its premises in the derivation below, each at its own standing.
The negativity of the Wigner marginals beyond two copies is the negativity of the swap-face square of the direct sum: the two objects carry the same sign at every one of the five copy counts from $N$ one to the parent bound three and the same first negative $N$, three, the total of the signed distribution equals the swap-face square plus the closure read at all four webs, and a second route from the many-body links reproduces the parent grids cell by cell, the first $N$ at which the two objects part reached not.
\end{chconditional}
\chstanding{a conditional theorem}
\chdeps{Definitions 1.1 to 1.6, Theorems 18.1 to 18.4 and Open results 18.5 and 18.6, Theorem 42.1 and Theorem 43.1}{Theorem 47.1}
One open result stands here. The Wigner function stands as the swap-face cell at four of five checks: the one-body record distribution over the pairs of a site and a direction at the two-dock electron at the standing placement (the position of a tensor in a network) is printed at sixteen cells and 3685 records with every cell nonnegative and both marginals equal to the washed readouts; the direct sum of the many-body result at the four tabulated joined webs decomposes into such cells with every cross pair between copies at the generic swap sign minus one, the pair part equal to the swap-face square and the closure part to the closure read; the marginals are negative at $N$ three and not at $N$ two, the negativity not in the joint cells alone; but the overlap of two distributions parts from the record overlap at the direct sum at both pairs tested, the electron with itself and with the muon body, so the entry of the Wigner function types a gloss at its fifth step.
\begin{chbridge}
The Wigner function of a state is a real quasi-probability distribution over phase space whose marginals are the position and the momentum distributions, for a finite system of dimension $d$ a function over the $d$ by $d$ phase-space cells that may be negative, and the overlap of two states is the phase-space integral of the product of their Wigner functions \cite{Wigner1932}; the exchange of two identical bosons or fermions enters the joint distribution of $N$ copies through the interference terms with the sign of the statistics \cite{Pauli1940}. The correspondence is by name: the compendium supplies a web of this theory read at its sinks, a doublet swap sign at the odd cycle, a Dirac form at an established generator and a finite record sum carrying the path-integral name at chapter~\ref{mod:II-1-state-space-from-the-act-outputs}, chapter~\ref{mod:II-5-indistinguishability-and-the-spin-statistics-sign}, chapter~\ref{mod:II-7-path-integral-as-the-finite-record-sum}, chapter~\ref{mod:III-5-spin-and-helicity}, the chapter titled The Dirac equation at the dock (chapter~\ref{mod:V-9-dirac-equation-at-the-dock}, among the instruments), chapter~\ref{mod:VI-2-w-and-the-z-and-their-masses}, chapter~\ref{mod:VI-2b-neutral-lepton-the-hub-unit-with-its-legs-open}, chapter~\ref{mod:VII-5-proton-to-electron-ratio-a-candidate}, chapter~\ref{mod:VIII-1-mass-unit-and-the-first-power}, chapter~\ref{mod:XIV-2-walls-census-in-physics-terms} and chapter~\ref{mod:XIV-7-what-remains-the-ungraded-page}, and the electron as the spin-one-half quantum of a field with its mass and spin as Lorentz labels and the muon as the spin-one-half quantum of a field with its mass and spin as Lorentz labels correspond to these objects by name.
\end{chbridge}
\begin{chproof}[exact part, finite part and census part]
The premises of this derivation are the axiom and the six definitions of chapter~\ref{mod:FRAMEWORK}, the theorem of chapter~\ref{mod:I-3-records-reads-and-the-born-count-the-gain-lemma}, the theorem of chapter~\ref{mod:II-4-measurement-the-wash-at-one-act-and-the-povm} and the theorem of chapter~\ref{mod:II-5-indistinguishability-and-the-spin-statistics-sign}. The chain of premises runs back to the axiom through those chapters, and the longest such chain in this compendium is twelve chapters long.
The signed distribution of the direct sum is built at the four tabulated joined webs from the grids of the parent result and from the many-body links at a second route, the Wigner negativity and the swap-face sign are compared side by side at every $N$ from one to the parent bound, the first negative $N$ is found at both objects, and the total of the signed distribution is compared with the swap-face square plus the closure read at every web; the theorem is conditional on the Wigner function as the swap-reading cell, an open result.
\textit{The exact part.}
The smallest joined web of two copies at which the swap-face square is read is the double pair, two copies of the two-dock body of chapter~\ref{mod:II-4-measurement-the-wash-at-one-act-and-the-povm} paired at both docks by two pairing links of type $(2,2)$, no link struck, the dressing nine. From the link list of that chapter, by the first-power rule stated there (the dressing multiplied by the sum of the measures of the unstruck links), we obtain the first powers of the three blocks, each copy with both dock links unstruck and the two pairing links together,
\begin{equation}
S_{1}(c_{1}) = S_{1}(c_{2}) = 9 \cdot (67 + 6 + 6) = 711, \qquad S_{1}(\text{pairing}) = 9 \cdot (4 + 4) = 72, \qquad S_{1} = 1494 ,
\end{equation}
and from the first powers, by the pair rule of that chapter, we obtain the pair sum as the square of the first power and the swap-face square as the washed square (the sum of the block squares) plus the cross terms with the pairing links minus the cross terms between the copies, every cross term between two copies carrying the sign minus one,
\begin{equation} \adjustbox{max width=\dimexpr\linewidth-3em\relax}{$\displaystyle \begin{gathered} P_{2} = 1494^{2} = 2232036, \\ Q_{\mathrm{swap}} = (711^{2} + 711^{2} + 72^{2}) + 2 \cdot 2 \cdot 72 \cdot 711 - 2 \cdot 711 \cdot 711 = 1016226 + 204768 - 1011042 = 209952 . \end{gathered}$} \end{equation}
The closure read of the direct sum at that web, the two-exchange part summed over its two hundred and fifty-four closure records, is an exact element of the ring in the basis of one, root two, root three and root six, taken here as a cell of the direct-sum tables of this result and re-computed on the page not,
\begin{equation}
R_{\mathrm{cl}} = \tfrac{25996610839}{2949120000} + \tfrac{1286}{225}\sqrt{2} + \tfrac{5977541}{147456000}\sqrt{3} + 0\sqrt{6} .
\end{equation}
From the signed distribution of that web over the pairs of a site and a direction, by summation, we obtain its total equal to the swap-face square plus the closure read, the two parts of the direct sum, and its closure marginals nonnegative at every one of the twenty-two sites, the sign of the direct sum at the swap nonnegative. At the pair web of two copies (one pairing link, the unpaired dock links struck, the blocks 657, 657 and 36 of that chapter) the same rule gives the swap-face square
\begin{equation}
Q_{\mathrm{swap}}(\text{pair}) = (657^{2} + 657^{2} + 36^{2}) + 2 \cdot 2 \cdot 36 \cdot 657 - 2 \cdot 657 \cdot 657 = 95904, \qquad \sum_{\text{site},\,\text{dir}} W = Q_{\mathrm{swap}} + R_{\mathrm{cl}} .
\end{equation}
From the marginals at every copy count from one to the parent bound three, by the sign of the first negative cell, we obtain the first negative copy count three at the Wigner marginals and three at the swap-face square, the same at both objects, and no parting at any copy count. A reader reproduces the sign at three copies from the same rule, at the chain of three copies (the blocks 2133, 1971 and 1971 at the dressing twenty-seven, the two pairing links 216 together) the swap-face square reads minus 9595827 and at the triangle (the blocks 2133, 2133 and 2133, the three pairing links 324 together) minus 9397539, negative at both where at two copies it is positive.
\begin{equation}
N_{\mathrm{neg}}(\text{Wigner}) = 3 = N_{\mathrm{neg}}(\text{swap}), \qquad N_{\mathrm{part}} = \text{none} .
\end{equation}
\textit{The finite part.}
The procedure. For each of the four tabulated joined webs (the pair and the double pair at two copies, the chain and the triangle at three, with the chain and the ring at four copies at the second route), build the signed distribution of the direct sum over the pairs of a site and a direction from the grids of the parent result, and rebuild it from the many-body links at a second route as exact elements of the ring. At every copy count compare the sign of the first negative marginal of the Wigner object with the sign of the swap-face square, find the first negative copy count at both, and check that the total of the signed distribution equals the swap-face square plus the closure read at every web. The smallest instance, the double pair at two copies, is worked above.
\textit{The census part.}
By exhaustive enumeration of the cells of the signed distribution at the four tabulated joined webs (eighteen cells at the double pair, every cell an exact element of the ring), the two routes agree cell by cell, the pair totals equal the pair sums at all six comparisons, and the negativity beyond two copies of the Wigner marginals is the negativity of the swap-face square at every copy count tested, the first negative copy count three at both.
The counter-tests alter the objects singly and read what moves. Striking the sign moves the first negative copy count at the main route while the two objects at the plus sign read no negative count, one cell moved. The three counter-tests stand at their pinned letters and the twelve controls read the same at both engines. The result was reproduced by a second engine built apart from the first, the two runs equal at every one of the cells of both with no differing key, and by the second route from the links. The closure read is a cell of the direct-sum tables of this result and is re-computed on the page not; every other quantity of the smallest instance follows from the link list and the two rules by a short program.
The machine-checked package named at the reproducibility appendix certifies the statements of this chapter under the definitions of this compendium and proves no theorem of the literature.
\chcert{6}
\end{chproof}
\begin{figure}[h]\centering
\setlength{\unitlength}{1pt}
\begin{picture}(260,150)
\qbezier(123.6,121.6)(103.5,101.5)(83.4,81.4)
\put(109.2,95.8){\makebox(0,0){\scriptsize $(2,2)$}}
\qbezier(136.4,28.4)(156.5,48.5)(176.6,68.6)
\put(150.8,54.2){\makebox(0,0){\scriptsize $(2,2)$}}
\qbezier(83.4,68.6)(103.5,48.5)(123.6,28.4)
\put(109.2,54.2){\makebox(0,0){\scriptsize $(2,2)$}}
\put(130.0,128.0){\circle{14}}
\put(130.0,128.0){\makebox(0,0){\footnotesize $2$}}
\put(130.0,138.0){\makebox(0,0)[b]{\scriptsize $A$}}
\put(77.0,75.0){\circle{14}}
\put(77.0,75.0){\makebox(0,0){\footnotesize $2$}}
\put(77.0,85.0){\makebox(0,0)[b]{\scriptsize $B$}}
\put(130.0,22.0){\circle{14}}
\put(130.0,22.0){\makebox(0,0){\footnotesize $2$}}
\put(130.0,32.0){\makebox(0,0)[b]{\scriptsize $A'$}}
\put(183.0,75.0){\circle{14}}
\put(183.0,75.0){\makebox(0,0){\footnotesize $2$}}
\put(183.0,85.0){\makebox(0,0)[b]{\scriptsize $B'$}}
\end{picture}
\caption{The smallest joined web of this chapter, two copies of a pair joined at a shared end into the pair web of two copies, every site a qubit. The signed distribution of its direct sum is summed over the pairs of a site and a direction.}
\end{figure}
\par\medskip\noindent{\itshape The pieces of quantum mechanics assembled, the next chapter gathers its corollaries at their roots.}\label{narr:close:2-8}\par
\section{The QM corollaries: the fifteen entries}\label{mod:II-9-qm-corollaries-the-fifteen-entries}
\chdesig{3-9}{2-9}{II-9-qm-corollaries-the-fifteen-entries}{II-9-D539176-w-tprd}
\label{narr:open:2-9}Theorem~47.1 places the quantum-mechanics corollaries at their established roots, so that the reader can see which familiar statement rests on which theorem of this Part.
\begin{chconditional}[The quantum-mechanics corollaries at their established roots]\label{box:II-9-qm-corollaries-the-fifteen-entries:1}
This result is conditional on the axiom text and on the results of the chapters named as its premises in the derivation below, each at its own standing.
The twenty-one quantum-mechanics entries of the remainder page of the glossary are each mapped to an established root and typed from the kind of the map and the standing of the root: six derived, a corollary at a theorem root with its sentence, its proof sketch and an exact check (the Wigner and Wigner-Eckart theorems at the rotation group of the derived directions, the optical theorem at unitarity from the flat pair, the projection postulate and decoherence at measurement as the wash at one act, the role of an end at the transport status of the end); twelve glosses, a continuum reading of an established object with no derivation claimed (among them the Lindblad form, the Zeno effect and the generalized measurement at the wash result, the Bell inequality and the Clebsch-Gordan coupling, the oscillator ladder and the coherent state at the first rungs, the classical limit and the Wentzel-Kramers-Brillouin expansion at the step generator); and three open at a named gate (the adiabatic theorem and the Berry phase at the ensemble chapters, the doubled site as a premise of the generic doublet swap); the nine exact checks hold as integer arithmetic, among them the Clauser-Horne-Shimony-Holt expectation at twice the classical square, and the twenty-one names stand as correspondences by name.
\end{chconditional}
\chstanding{a conditional theorem}
\chdeps{Definitions 1.1 to 1.6, Theorems 28.1 and 28.2, Open results 28.3 and 28.4 and Theorem 28.5, Theorem 35.1, Theorem 40.1, Theorem 42.1, Theorem 43.1, Theorem 44.1 and Theorem 46.1}{Chapter 3-9 and Open result 228.1}
One open result stands here. Four corollaries re-derive from named clauses at three of four, the Bell inequality at the rotated settings of the derived rotation group with the Clauser-Horne-Shimony-Holt maximum twice the square root of two, the Clebsch-Gordan coupling at the two ends of a link, and the tunnelling amplitude at the dark (3,3) links at order two, their twelve exact checks holding; the Wigner function is the first naming, its missing clause the swap face as a cell, the record distribution at the electron web of the construction carrying no negative cell while the swap face of the many-body result is negative beyond two copies.
\begin{chbridge}
The textbook corollaries of quantum mechanics beyond its axioms, from the Wigner and Wigner-Eckart theorems and the optical theorem through the Bell inequality, the Clebsch-Gordan coupling, the tunnelling amplitude, the Wigner function, the Lindblad form and the Zeno effect to the classical limit, the adiabatic theorem and the Berry phase, are each a derivation from the axioms at its standard setting \cite{Sakurai1994,Nielsen2000}. The correspondence is by name.
\end{chbridge}
\begin{chproof}[exact part and census part]
The premises of this derivation are the axiom and the six definitions of chapter~\ref{mod:FRAMEWORK}, the theorem of chapter~\ref{mod:I-6-read-at-a-register-the-probe-and-the-measurement}, the theorem of chapter~\ref{mod:II-1-state-space-from-the-act-outputs}, the theorem of chapter~\ref{mod:II-2-born-rule-and-unitarity}, the theorem of chapter~\ref{mod:II-4-measurement-the-wash-at-one-act-and-the-povm}, the theorem of chapter~\ref{mod:II-5-indistinguishability-and-the-spin-statistics-sign}, the theorem of chapter~\ref{mod:II-6-energy-and-the-hamiltonian-the-first-power-operator-and-the-step} and the theorem of chapter~\ref{mod:II-8-wigner-function-and-the-many-body-sign}. The chain of premises runs back to the axiom through those chapters, and the longest such chain in this compendium is twelve chapters long.
The twenty-one entries are located at their roots, every root read at its standing, the kind of every map typed from its sketch as a corollary with an exact check, a naming or a gate, and the nine exact checks run as integer arithmetic at the construction or at the one link, the deciding reading at the gate of the page with the page-alone reading printed beside; the theorem is conditional on its roots and on the open results it rests on, among them the first rungs as the iterated closure and the energy as the step generator, and on the glossary of the textbook corpus.
\textit{The exact part.}
From the three derived directions at the qubit end of the read link, by the signed permutation maps of the three axes, we obtain forty-eight maps falling in exactly two kinds by determinant, twenty-four proper and twenty-four improper, every one preserving the Born pairing of the six axis states at all one thousand seven hundred and twenty-eight pairs, the exact check of the first corollary.
\begin{equation}
\#\{\text{maps}\} = 48 = 24 + 24, \qquad \det = \pm 1, \qquad \#\{\text{Born pairs preserved}\} = 1728 \ \text{among } 1728 .
\end{equation}
\textit{The census part.}
By exhaustive enumeration of the twenty-one quantum-mechanics entries of the remainder page of the glossary (every entry read at its corpus name and located at its root, a result of this compendium read at its standing, fourteen roots theorems, five roots open results, two entries gated with no root), the kind of every map is typed from its sketch as a corollary with an exact check, a gloss or a gate, six derived corollaries, each with its sentence, its proof sketch and an exact check run as integer arithmetic at the construction or at the one link, twelve glosses with no derivation claimed, and three gated entries. The nine exact checks hold. The smallest case worked here is the first corollary, the projection postulate, it is the wash at one act, its root the measurement result, the far end of the pair read at a basis vector and washed leaving the near end in the matching projector, and its check the signed permutation maps of the three derived directions at the qubit, which fall in exactly two kinds by determinant (twenty-four proper and twenty-four improper among the forty-eight maps) with the Born pairing of the six axis states preserved at all 1728 pairs.
The counter-tests alter the objects singly and read what moves. Striking one root moves five entries from derived or gloss to open, all five of them, the projection postulate the first. The three counter-tests stand at their pinned letters and the ten controls read the same at both engines, among them the twenty-one names equal to their corpus rows. The result was reproduced by a second engine built apart from the first, the two runs equal at every one of the cells of both.
The machine-checked package named at the reproducibility appendix certifies the statements of this chapter under the definitions of this compendium and proves no theorem of the literature.
\chcert{21}
\end{chproof}
\begin{figure}[h]\centering
\setlength{\unitlength}{1pt}
\begin{picture}(260,150)
\qbezier(86.0,75.0)(130.0,75.0)(174.0,75.0)
\put(130.0,83.0){\makebox(0,0){\scriptsize $(2,3)$}}
\put(77.0,75.0){\circle{14}}
\put(77.0,75.0){\makebox(0,0){\footnotesize $3$}}
\put(77.0,85.0){\makebox(0,0)[b]{\scriptsize $q$}}
\put(183.0,75.0){\circle{14}}
\put(183.0,75.0){\makebox(0,0){\footnotesize $2$}}
\put(183.0,85.0){\makebox(0,0)[b]{\scriptsize $b$}}
\end{picture}
\caption{The read link at its qubit end, the object at which the exact checks of the derived corollaries are run, the three derived directions of the qubit end and the six axis states of their Born pairing.}
\end{figure}
\section{The three premises as corollaries of the read cap}\label{mod:II-9-D539176-w-tprd}
The information a measurement can extract from a quantum system is capped by the dimension of the system: a d-dimensional system carries at most log d bits at a single read (the Holevo bound), a projective measurement of a maximally entangled pair leaves the partner in a definite state (the collapse postulate), and the composition of two systems is their tensor product; these three statements are postulates of the standard formalism and are not derived within it \cite{Holevo1973,Nielsen2000}. The quantum-mechanical corollaries are gathered at section~\ref{mod:II-9-qm-corollaries-the-fifteen-entries} and the axiom and the six definitions at section~\ref{mod:I-1-axiom-and-the-six-definitions}.
The three premises that the invariance census types as premises, that an end is open, that a saturation stops and that a join passes a read, are posed as corollaries of the read cap. At every one of the 1262 ends of the eight webs of the cap census, the open set of the standing definition at its four readings is compared with the set of ends the read leaves open at the far end and with the set of ends the wash still moves on the state: no reading meets either set at every end, the first premise parts, the second parts likewise and the third is reproduced at the read cap.
The standing open set at its principal reading agrees with the read-open set at 500 of the 1262 ends and with the wash-moved set at 881, the first parting at an end of the universe web that the definition marks open while the read reads it closed and the wash leaves it unmoved. Beside the four readings, the read-open set at the near end equals the standing open set taken at the near end, and the wash-moved set at the end register equals the read-open set at 1262 of 1262 ends: the read cap is the wash fixed point on the state, at the larger end of every unequal link, the end the standing definition marks washed. The stop condition agrees with read-closed-everywhere at 2 of the 8 webs, the first parting a web whose saturation stops at wave 0 while no wave is read closed everywhere; the join by docking agrees with the pairing move passing a read at the read cap of its (2,2) contraction, read-closed at both ends.
The result is an open one, an open result: the premises are not re-derived from the read cap at any of the four readings of the standing definition, and the result stands at the weaker reading with the first premise, open, named; no comparison with observation is made here.
What stays open is the premise itself, standing as a word rather than a corollary; the definition of open at the larger end (section~\ref{mod:I-1-D539391-w-dole}) is posed from this identity as its mirror, and the premise census of the mass chain (section~\ref{mod:V-1-D537668-the-mass-chain-premise-discharge-census}) keeps its typing.
\par\medskip\noindent{\itshape One subject of the Part remains, the system that is not closed.}\label{narr:close:2-9}\par
\section{The instrument tables of this Part}\label{mod:III-L-the-instrument-tables-at-the-repository}
\chdesig{3-L}{}{III-L-the-instrument-tables-at-the-repository}{III-L-the-instrument-tables-at-the-repository}
\begin{chtable}[The instrument tables of this Part]\label{box:III-L-the-instrument-tables-at-the-repository:1}
This Part lists 2 instrument chapters, printed among the instruments at the end of this compendium as chapters~\ref{mod:II-10-decoherence-and-the-open-system} to~\ref{mod:II-10-D536583-psyo-t1}, and the table beneath names each one with its key cells.
\end{chtable}
\chstanding{instrument}
\chdeps{Definitions 1.1 to 1.6}{no theorem beyond it at this release}
The listing moves no cell and states no result of its own. Each instrument stands at the repository as its chapter stood at the previous release of this compendium, and its tables are read there.
\begin{chbridge}
The tables listed here are instruments of this compendium, each listed under the title of its chapter.
\end{chbridge}
\begin{chproof}[instrument]
The listing in full. For each instrument chapter, printed among the instruments at the end of this compendium, the table gives its title, the result boxes it carries, the standing of each box with its citation, and the count of objects its certificate sentence names where it prints one.
{\small
\begin{longtable}{p{38mm}p{63mm}p{30mm}p{14mm}}
\hline
chapter & results & standing & objects \\
\hline
\endhead
Decoherence and the open system & The chapter counted empty & instrument & -- \\
Decoherence and the open system & The open system as the wash at a sub-web & theorem & 1755 \\
\hline
\end{longtable}
}
The rule that produced it. Every chapter of this Part whose first standing beyond a placeholder reads instrument is printed among the instruments at the end of this compendium with its tables, and its key cells are the titles of its boxes, its standing lines and its certificate count, read from its text as it stood at the previous release.
No derivation stands in this chapter; the tables stand at the repository.
\emph{Check.} Every file listed was re-hashed at the repository and read the same as its recorded hash, and the listing was re-cut from the files by a second route.
Nothing stands open at this listing. At this release one result of this Part is re-graded at its chapter: the path integral as the finite record sum (the chapter 3-7), an open result re-posed at its first unkeyed clause.
\end{chproof}
\label{chend:III-L-the-instrument-tables-at-the-repository}
